% ==============================================================================
% Softwareentwurf – Vintage Car Garage
% Benötigte Pakete im Hauptdokument:
%   \usepackage{longtable}
%   \usepackage{colortbl}
%   \usepackage{enumitem}
%   \usepackage{amssymb}
%   \usepackage{listings}
%   \usepackage[utf8]{inputenc}
%   \definecolor{tableHeader}{RGB}{52,73,94}
%   \newcolumntype{L}{>{\raggedright\arraybackslash}p{0.28\linewidth}}
%   \newcolumntype{R}{>{\raggedright\arraybackslash}p{0.65\linewidth}}
% ==============================================================================


% ==============================================================================
% TEIL 0: NEU AUFGENOMMENE VERWALTUNGSOPERATIONEN – USE-CASE-ZUORDNUNG
% ==============================================================================
\phantomsection
\addcontentsline{toc}{subsection}{Use-Case-Zuordnung der ergänzten Operationen}
\subsection*{Use-Case-Zuordnung der erg\"{a}nzten Operationen}

Die folgenden – zuvor auskommentierten – Verwaltungsoperationen wurden in das
erste Release aufgenommen. Sie verfeinern bestehende Use Cases der beiden Akteure
(keine neuen Use Cases). Jede Operation besitzt einen operationellen Vertrag
(TEIL 2) und ein Entwurfssequenzdiagramm (TEIL 4).

\begin{longtable}{|l|l|l|p{0.34\linewidth}|}
\hline
\rowcolor{tableHeader}
\textbf{Operation} & \textbf{Vertrag} & \textbf{Akteur} & \textbf{Use Case} \\ \hline
\endhead
getAllRequests & CO\_UC\_E2\_2 & Einlagerer & UC\_E2 – Stellplatzangebot anfordern \\ \hline
deleteRequest & CO\_UC\_E2\_3 & Einlagerer & UC\_E2 – Stellplatzangebot anfordern \\ \hline
removeBookedService & CO\_UC\_E6\_2 & Einlagerer & UC\_E6 – Optionale Lagerservices buchen \\ \hline
getAllBookedAppointments & CO\_UC\_E5\_4 & Einlagerer & UC\_E5 – Ein-/Auslagerungstermin koordinieren \\ \hline
cancelAppointment & CO\_UC\_E5\_5 & Einlagerer & UC\_E5 – Ein-/Auslagerungstermin koordinieren \\ \hline
getAllAppointments & CO\_UC\_E5\_6 & Lagerhalter & UC\_E5 – Ein-/Auslagerungstermin koordinieren \\ \hline
deleteAppointment & CO\_UC\_E5\_7 & Lagerhalter & UC\_E5 – Ein-/Auslagerungstermin koordinieren \\ \hline
getAllOffers & CO\_UC\_L7\_3 & Lagerhalter & UC\_L7 – Stellplatzanfragen verwalten \\ \hline
deleteOffer & CO\_UC\_L7\_4 & Lagerhalter & UC\_L7 – Stellplatzanfragen verwalten \\ \hline
\end{longtable}

Zus\"{a}tzlich wurde der bereits entworfene Endpunkt \texttt{getAvailableAppointments}
(CO\_UC\_E5\_2, Einlagerer, UC\_E5) reaktiviert und auf die Signatur
\texttt{(carOwnerId, storageOwnerId)} gebracht.
\newpage

% ==============================================================================
% TEIL 1: SYSTEMSEQUENZDIAGRAMME (ANNOTIERT)
% ==============================================================================

% ------------------------------------------------------------------------------
% UC_B1.1.1 – Lagerhalter registrieren
% ------------------------------------------------------------------------------
\phantomsection
\addcontentsline{toc}{subsubsection}{UC\_B1.1.1 – Lagerhalter registrieren}
\subsubsection*{UC\_B1.1.1 – Lagerhalter registrieren}

\begin{lstlisting}[basicstyle=\ttfamily\small, breaklines=true]
sequenceDiagram
    actor Lagerhalter
    participant System

    Lagerhalter->>System: 1. Ruft die Ansicht "Registrieren" auf
    Lagerhalter->>System: 2. wählt die Stakeholder-Rolle "Lagerhalter" aus
    Lagerhalter->>System: 3. Fuellt Profildaten (EMailAdresse, Passwort, Firmenname,
                              Standort, Lagerbedingungen, Kapazität, Kosten,
                              Servicezeit, Markenspezialisierung (opt.)) aus
    Lagerhalter->>System: 4. Bestätigt die Eingaben
    note over System: Vertrag CO_UC_B1_1: createStorageOwner
    System->>System: 5. Prueft EMailAdresse & Passwort auf Gueltigkeit

    alt 5.1 Ausnahmeszenario: E-Mail-Adresse bereits registriert
        System-->>Lagerhalter: 5.1.1 Informiert den Benutzer ueber das bereits
                                      existierende Konto
        System-->>Lagerhalter: 5.1.2 Bittet den Benutzer, sich anzumelden
                                      oder eine andere E-Mail zu nutzen
    else 5.2 Ausnahmeszenario: Passwoerter stimmen nicht ueberein
        System-->>Lagerhalter: 5.2.1 Informiert den Benutzer, dass die Passwoerter
                                      nicht uebereinstimmen
        System-->>Lagerhalter: 5.2.2 Bittet den Benutzer, die Eingabe zu korrigieren
    end

    System->>System: 6. Speichert die Daten dauerhaft
    System-->>Lagerhalter: 7. Bestätigt dem Benutzer die erfolgreiche Registrierung
\end{lstlisting}

\newpage

% ------------------------------------------------------------------------------
% UC_B1.1.2 – Einlagerer registrieren
% ------------------------------------------------------------------------------
\phantomsection
\addcontentsline{toc}{subsubsection}{UC\_B1.1.2 – Einlagerer registrieren}
\subsubsection*{UC\_B1.1.2 – Einlagerer registrieren}

\begin{lstlisting}[basicstyle=\ttfamily\small, breaklines=true]
sequenceDiagram
    actor Einlagerer
    participant System

    Einlagerer->>System: 1. Ruft die Ansicht "Registrieren" auf
    Einlagerer->>System: 2. wählt die Stakeholder-Rolle "Einlagerer" aus
    Einlagerer->>System: 3. Fuellt Profildaten (EMailAdresse, Passwort,
                              Vorname, Nachname) aus
    Einlagerer->>System: 4. Bestätigt die Eingaben
    note over System: Vertrag CO_UC_B1_2: createCarOwner
    System->>System: 5. Prueft EMailAdresse & Passwort auf Gueltigkeit

    alt 5.1 Ausnahmeszenario: E-Mail-Adresse bereits registriert
        System-->>Einlagerer: 5.1.1 Informiert den Benutzer ueber das
                                     bereits existierende Konto
        System-->>Einlagerer: 5.1.2 Bittet den Benutzer, sich anzumelden
                                     oder eine andere E-Mail zu nutzen
    else 5.2 Ausnahmeszenario: Passwoerter stimmen nicht ueberein
        System-->>Einlagerer: 5.2.1 Informiert den Benutzer, dass die Passwoerter
                                     nicht uebereinstimmen
        System-->>Einlagerer: 5.2.2 Bittet den Benutzer, die Eingabe zu korrigieren
    end

    System->>System: 6. Speichert die Daten dauerhaft
    System-->>Einlagerer: 7. Bestätigt dem Benutzer die erfolgreiche Registrierung
\end{lstlisting}

\newpage

% ------------------------------------------------------------------------------
% UC_B2 – Benutzer anmelden
% ------------------------------------------------------------------------------
\phantomsection
\addcontentsline{toc}{subsubsection}{UC\_B2 – Benutzer anmelden}
\subsubsection*{UC\_B2 – Benutzer anmelden}

\begin{lstlisting}[basicstyle=\ttfamily\small, breaklines=true]
sequenceDiagram
    actor Benutzer
    participant System

    Benutzer->>System: 1. Ruft die Ansicht "Anmelden" auf
    Benutzer->>System: 2. Gibt Zugangsdaten (EMailAdresse, Passwort) ein
    Benutzer->>System: 3. Bestätigt die Eingaben
    note over System: Vertrag CO_UC_B2_1: authenticate
    System->>System: 4. Prueft die Zugangsdaten auf Gueltigkeit

    alt 4.1 Ausnahmeszenario: Zugangsdaten ungueltig
        System-->>Benutzer: 4.1.1 Informiert den Benutzer ueber fehlerhafte Zugangsdaten
        System-->>Benutzer: 4.1.2 Bittet den Benutzer, die Eingabe zu korrigieren
    else Hauptszenario: Zugangsdaten gueltig
        System->>System: 5. Meldet den Benutzer an
    end
\end{lstlisting}

\newpage

% ------------------------------------------------------------------------------
% UC_B3.1 – Lagerhalter verwalten
% ------------------------------------------------------------------------------
\phantomsection
\addcontentsline{toc}{subsubsection}{UC\_B3.1 – Lagerhalter verwalten}
\subsubsection*{UC\_B3.1 – Lagerhalter verwalten}

\begin{lstlisting}[basicstyle=\ttfamily\small, breaklines=true]
sequenceDiagram
    actor Lagerhalter
    participant System

    Lagerhalter->>System: 1. Ruft die Ansicht "Benutzerprofil" auf
    note over System: Vertrag CO_UC_B3_1: getStorageOwner
    System-->>Lagerhalter: 2. Zeigt die aktuellen Profildaten (EMailAdresse, Passwort,
                                Firmenname, Standort, Lagerbedingungen, Kapazität,
                                Kosten, Servicezeit, Markenspezialisierung (opt.)) an
    Lagerhalter->>System: 3. Passt seine Profildaten an
    Lagerhalter->>System: 4. Bestätigt die Eingaben
    note over System: Vertrag CO_UC_B3_2: updateStorageOwner
    System->>System: 5. Prueft EMailAdresse & Passwort auf Gueltigkeit

    alt 5.1 Ausnahmeszenario: E-Mail-Adresse bereits registriert
        System-->>Lagerhalter: 5.1.1 Informiert den Benutzer ueber den Konflikt
                                      und bittet um Korrektur
    else 5.2 Ausnahmeszenario: Passwoerter stimmen nicht ueberein
        System-->>Lagerhalter: 5.2.1 Informiert den Benutzer, dass die Passwoerter
                                      nicht uebereinstimmen
        System-->>Lagerhalter: 5.2.2 Bittet den Benutzer, die Eingabe zu korrigieren
    else Hauptszenario: Daten gueltig
        System->>System: 6. Speichert die Daten dauerhaft
        System-->>Lagerhalter: 7. Aktualisiert die Ansicht
    end
\end{lstlisting}

\newpage

% ------------------------------------------------------------------------------
% UC_B3.2 – Einlagerer verwalten
% ------------------------------------------------------------------------------
\phantomsection
\addcontentsline{toc}{subsubsection}{UC\_B3.2 – Einlagerer verwalten}
\subsubsection*{UC\_B3.2 – Einlagerer verwalten}

\begin{lstlisting}[basicstyle=\ttfamily\small, breaklines=true]
sequenceDiagram
    actor Einlagerer
    participant System

    Einlagerer->>System: 1. Ruft die Ansicht "Benutzerprofil" auf
    note over System: Vertrag CO_UC_B3_1: getCarOwner
    System-->>Einlagerer: 2. Zeigt die aktuellen Profildaten (EMailAdresse, Passwort,
                               Vorname, Nachname) an
    Einlagerer->>System: 3. Passt seine Profildaten an
    Einlagerer->>System: 4. Bestätigt die Eingaben
    note over System: Vertrag CO_UC_B3_2: updateCarOwner
    System->>System: 5. Prueft EMailAdresse & Passwort auf Gueltigkeit

    alt 5.1 Ausnahmeszenario: E-Mail-Adresse bereits registriert
        System-->>Einlagerer: 5.1.1 Informiert den Benutzer ueber den Konflikt
                                     und bittet um Korrektur
    else 5.2 Ausnahmeszenario: Passwoerter stimmen nicht ueberein
        System-->>Einlagerer: 5.2.1 Informiert den Benutzer, dass die Passwoerter
                                     nicht uebereinstimmen
        System-->>Einlagerer: 5.2.2 Bittet den Benutzer, die Eingabe zu korrigieren
    else Hauptszenario: Daten gueltig
        System->>System: 6. Speichert die Daten dauerhaft
        System-->>Einlagerer: 7. Aktualisiert die Ansicht
    end
\end{lstlisting}

\newpage

% ------------------------------------------------------------------------------
% UC_L1 – Serviceportfolio verwalten
% ------------------------------------------------------------------------------
\phantomsection
\addcontentsline{toc}{subsubsection}{UC\_L1 – Serviceportfolio verwalten}
\subsubsection*{UC\_L1 – Serviceportfolio verwalten}

\begin{lstlisting}[basicstyle=\ttfamily\small, breaklines=true]
sequenceDiagram
    actor Lagerhalter
    participant System

    Lagerhalter->>System: 1. Ruft die Ansicht "Serviceportfolio" auf
    note over System: Vertrag CO_UC_L1_1: getAllServices
    System-->>Lagerhalter: 2. Zeigt die Uebersicht der angebotenen Services an
    Lagerhalter->>System: 3. Passt das Serviceangebot (Titel, Beschreibung,
                              Preis, IstOptional) an
    note over System: Vertrag CO_UC_L1_2: createService
    note over System: Vertrag CO_UC_L1_3: updateService
    note over System: Vertrag CO_UC_L1_4: deleteService
    Lagerhalter->>System: 4. Bestätigt die Eingaben
    System->>System: 5. Speichert die Daten dauerhaft
    System-->>Lagerhalter: 6. Aktualisiert die Ansicht
\end{lstlisting}

\newpage

% ------------------------------------------------------------------------------
% UC_L2 – Stellplatzkapazitäten verwalten
% ------------------------------------------------------------------------------
\phantomsection
\addcontentsline{toc}{subsubsection}{UC\_L2 – Stellplatzkapazit\"{a}ten verwalten}
\subsubsection*{UC\_L2 – Stellplatzkapazit\"{a}ten verwalten}

\begin{lstlisting}[basicstyle=\ttfamily\small, breaklines=true]
sequenceDiagram
    actor Lagerhalter
    participant System

    Lagerhalter->>System: 1. Ruft die Ansicht "Stellplatzkapazitäten" auf
    note over System: Vertrag CO_UC_B3_1: getStorageOwner
    System-->>Lagerhalter: 2. Zeigt die Uebersicht der Stellplätze an
    Lagerhalter->>System: 3. Aktualisiert die Kapazitäten des Lagers
    Lagerhalter->>System: 4. Bestätigt die Eingaben
    note over System: Vertrag CO_UC_B3_2: updateStorageOwner
    System->>System: 5. Speichert die Daten dauerhaft
    System-->>Lagerhalter: 6. Aktualisiert die Ansicht
\end{lstlisting}

\newpage

% ------------------------------------------------------------------------------
% UC_L3 – Fahrzeugstatus aktualisieren
% ------------------------------------------------------------------------------
\phantomsection
\addcontentsline{toc}{subsubsection}{UC\_L3 – Fahrzeugstatus aktualisieren}
\subsubsection*{UC\_L3 – Fahrzeugstatus aktualisieren}

\begin{lstlisting}[basicstyle=\ttfamily\small, breaklines=true]
sequenceDiagram
    actor Lagerhalter
    participant System

    Lagerhalter->>System: 1. Ruft die Ansicht "Fahrzeuguebersicht" auf
    note over System: Vertrag CO_UC_L3_1: getStoredCars
    System-->>Lagerhalter: 2. Zeigt die Uebersicht der Fahrzeuge an
    Lagerhalter->>System: 3. Passt die Statusdaten (Fahrbereitschaft, Wartungsstand)
                              eines Fahrzeugs an
    Lagerhalter->>System: 4. Bestätigt die Eingaben
    note over System: Vertrag CO_UC_L3_2: updateStoredCarStatus
    System->>System: 5. Speichert die Daten dauerhaft
    System-->>Lagerhalter: 6. Aktualisiert die Ansicht
\end{lstlisting}

\newpage

% ------------------------------------------------------------------------------
% UC_L6 – Optionale Services verwalten
% ------------------------------------------------------------------------------
\phantomsection
\addcontentsline{toc}{subsubsection}{UC\_L6 – Optionale Services verwalten}
\subsubsection*{UC\_L6 – Optionale Services verwalten}

\begin{lstlisting}[basicstyle=\ttfamily\small, breaklines=true]
sequenceDiagram
    actor Lagerhalter
    participant System

    Lagerhalter->>System: 1. Ruft die Ansicht "Serviceportfolio" auf
    note over System: Vertrag CO_UC_L6_1: getOptionalServices
    System-->>Lagerhalter: 2. Zeigt die Uebersicht der Services an
    Lagerhalter->>System: 3. Passt das Serviceangebot (Titel, Beschreibung,
                              Preis, IstOptional) an
    note over System: Vertrag CO_UC_L1_2: createService
    note over System: Vertrag CO_UC_L1_3: updateService
    note over System: Vertrag CO_UC_L1_4: deleteService
    Lagerhalter->>System: 4. Bestätigt die Eingaben
    System->>System: 5. Speichert die Daten dauerhaft
    System-->>Lagerhalter: 6. Aktualisiert die Ansicht
\end{lstlisting}

\newpage

% ------------------------------------------------------------------------------
% UC_L7 – Stellplatzanfragen verwalten
% ------------------------------------------------------------------------------
\phantomsection
\addcontentsline{toc}{subsubsection}{UC\_L7 – Stellplatzanfragen verwalten}
\subsubsection*{UC\_L7 – Stellplatzanfragen verwalten}

\begin{lstlisting}[basicstyle=\ttfamily\small, breaklines=true]
sequenceDiagram
    actor Lagerhalter
    participant System
    actor Einlagerer

    Lagerhalter->>System: 1. Ruft die Ansicht "Stellplatzanfragen" auf
    note over System: Vertrag CO_UC_L7_1: getIncomingRequests
    System-->>Lagerhalter: 2. Zeigt die Uebersicht der offenen Anfragen an
    Lagerhalter->>System: 3. wählt eine Anfrage aus
    Lagerhalter->>System: 4. Erstellt ein Angebot (Beschreibung, Kosten)
    note over System: Vertrag CO_UC_L7_2: createOffer
    System->>System: 5. Initialisiert Angebotstatus
    Lagerhalter->>System: 6. Bestätigt die Eingaben
    System->>System: 7. Speichert die Daten dauerhaft
    System-->>Einlagerer: 8. Uebermittelt das Angebot an den Einlagerer
    Lagerhalter->>System: 9. Ruft die Uebersicht seiner Angebote auf
    note over System: Vertrag CO_UC_L7_3: getAllOffers
    Lagerhalter->>System: 10. Zieht ein Angebot zurueck
    note over System: Vertrag CO_UC_L7_4: deleteOffer
\end{lstlisting}

\newpage

% ------------------------------------------------------------------------------
% UC_E1 – Stellplatz suchen
% ------------------------------------------------------------------------------
\phantomsection
\addcontentsline{toc}{subsubsection}{UC\_E1 – Stellplatz suchen}
\subsubsection*{UC\_E1 – Stellplatz suchen}

\begin{lstlisting}[basicstyle=\ttfamily\small, breaklines=true]
sequenceDiagram
    actor Einlagerer
    participant System

    Einlagerer->>System: 1. Ruft die Ansicht "Stellplatzsuche" auf
    Einlagerer->>System: 2. Sucht nach Standort (Stadt) und bestätigt die Eingaben
    note over System: Vertrag CO_UC_E1_1: searchByCity
    System-->>Einlagerer: 3. Zeigt die Uebersicht der Suchergebnisse (Standort) an

    alt 3.1 Ausnahmeszenario: Keine Treffer
        System-->>Einlagerer: 3.1.1 Informiert den Einlagerer ueber die leere Trefferliste
    end
\end{lstlisting}

\newpage

% ------------------------------------------------------------------------------
% UC_E2 – Stellplatzangebot anfordern
% ------------------------------------------------------------------------------
\phantomsection
\addcontentsline{toc}{subsubsection}{UC\_E2 – Stellplatzangebot anfordern}
\subsubsection*{UC\_E2 – Stellplatzangebot anfordern}

\begin{lstlisting}[basicstyle=\ttfamily\small, breaklines=true]
sequenceDiagram
    actor Einlagerer
    participant System
    actor Lagerhalter

    Einlagerer->>System: 1. Ruft die Ansicht "Lagerstandort" auf
    note over System: Vertrag CO_UC_E1_1: searchByCity
    System-->>Einlagerer: 2. Zeigt die Standortdetails (Strasse, Hausnummer,
                               PLZ, Stadt) an
    Einlagerer->>System: 3. Formuliert seine Anfrage (Typ, Beschreibung)
    Einlagerer->>System: 4. Bestätigt die Eingaben
    note over System: Vertrag CO_UC_E2_1: createRequest
    System->>System: 5. Speichert die Daten dauerhaft
    System-->>Lagerhalter: 6. Uebermittelt die Anfrage an den Lagerhalter
    Einlagerer->>System: 7. Ruft die Uebersicht seiner Anfragen auf
    note over System: Vertrag CO_UC_E2_2: getAllRequests
    Einlagerer->>System: 8. Zieht eine Anfrage zurueck
    note over System: Vertrag CO_UC_E2_3: deleteRequest
\end{lstlisting}

\newpage

% ------------------------------------------------------------------------------
% UC_E3 – Stellplatzangebote verwalten
% ------------------------------------------------------------------------------
\phantomsection
\addcontentsline{toc}{subsubsection}{UC\_E3 – Stellplatzangebote verwalten}
\subsubsection*{UC\_E3 – Stellplatzangebote verwalten}

\begin{lstlisting}[basicstyle=\ttfamily\small, breaklines=true]
sequenceDiagram
    actor Einlagerer
    participant System
    actor Lagerhalter

    Einlagerer->>System: 1. Ruft die Ansicht "Stellplatzangebote" auf
    note over System: Vertrag CO_UC_E3_1: getReceivedOffers
    System-->>Einlagerer: 2. Zeigt die Uebersicht der offenen Angebote an
    Einlagerer->>System: 3. wählt ein Angebot aus
    System-->>Einlagerer: 4. Zeigt die Details des Angebots (Beschreibung, Kosten) an

    alt 5.1 Ausnahmeszenario: Angebot wird abgelehnt
        Einlagerer->>System: 5.1.1 wählt die Ablehnung des Angebots aus
        note over System: Vertrag CO_UC_E3_2: respondToOffer (accepted=false)
        System->>System: 5.1.2 Aktualisiert den Angebotsstatus
        System->>System: 5.1.3 Speichert die Daten dauerhaft
        System-->>Lagerhalter: 5.1.4 Informiert den Lagerhalter
    else Hauptszenario: Angebot wird angenommen
        Einlagerer->>System: 5. Nimmt das Angebot an
        note over System: Vertrag CO_UC_E3_2: respondToOffer (accepted=true)
        System->>System: 6. Speichert die Daten dauerhaft
        System-->>Lagerhalter: 7. Informiert den Lagerhalter
    end
\end{lstlisting}

\newpage

% ------------------------------------------------------------------------------
% UC_E5 – Ein- & Auslagerungstermin koordinieren
% ------------------------------------------------------------------------------
\phantomsection
\addcontentsline{toc}{subsubsection}{UC\_E5 – Ein- \& Auslagerungstermin koordinieren}
\subsubsection*{UC\_E5 – Ein- \& Auslagerungstermin koordinieren}

\begin{lstlisting}[basicstyle=\ttfamily\small, breaklines=true]
sequenceDiagram
    actor Lagerhalter
    actor Einlagerer
    participant System

    Lagerhalter->>System: 0. Hat zuvor verfuegbare Termine eingetragen
    note over System: Vertrag CO_UC_E5_1: createAppointment
    note over System: Vertrag CO_UC_E5_6: getAllAppointments
    note over System: Vertrag CO_UC_E5_7: deleteAppointment

    Einlagerer->>System: 1. Ruft die Ansicht "Lagerhalterkalender" auf
    note over System: Vertrag CO_UC_E5_2: getAvailableAppointments
    Einlagerer->>System: 2. wählt den Betreff fuer den Termin aus
    System-->>Einlagerer: 3. Zeigt die verfuegbaren Termine des Lagerhalters an

    alt 3.1 Ausnahmeszenario: Keine freien Termine verfuegbar
        System-->>Einlagerer: 3.1.1 Informiert den Einlagerer darueber,
                                     dass keine Termine verfuegbar sind
    else Hauptszenario: Termine verfuegbar
        Einlagerer->>System: 4. wählt einen passenden Termin (Datum, Uhrzeit) aus
        Einlagerer->>System: 5. Bestätigt die Eingaben
        note over System: Vertrag CO_UC_E5_3: bookAppointment
        System->>System: 6. Prueft die Daten auf Konflikte
        System->>System: 7. Speichert die Daten dauerhaft
        System->>System: 8. Aktualisiert den Fahrzeugstatus (Fahrbereitschaft)
    end
    Einlagerer->>System: 9. Ruft die Uebersicht seiner gebuchten Termine auf
    note over System: Vertrag CO_UC_E5_4: getAllBookedAppointments
    Einlagerer->>System: 10. Storniert einen gebuchten Termin
    note over System: Vertrag CO_UC_E5_5: cancelAppointment
\end{lstlisting}

\newpage

% ------------------------------------------------------------------------------
% UC_E6 – Optionale Lagerservices buchen
% ------------------------------------------------------------------------------
\phantomsection
\addcontentsline{toc}{subsubsection}{UC\_E6 – Optionale Lagerservices buchen}
\subsubsection*{UC\_E6 – Optionale Lagerservices buchen}

\begin{lstlisting}[basicstyle=\ttfamily\small, breaklines=true]
sequenceDiagram
    actor Lagerhalter
    actor Einlagerer
    participant System

    Lagerhalter->>System: 0. Hat zuvor optionale Services angelegt
    note over System: Vertrag CO_UC_L1_2: createService

    Einlagerer->>System: 1. Ruft die Ansicht "Serviceportfolio" auf
    note over System: Vertrag CO_UC_L6_1: getOptionalServices
    System-->>Einlagerer: 2. Zeigt die Uebersicht der buchbaren Services an

    alt Service buchen
        Einlagerer->>System: 3a. wählt die gewuenschten Services aus
        Einlagerer->>System: 4a. Bestätigt die Auswahl
        note over System: Vertrag CO_UC_E6_1: bookOptionalService
        System->>System: 5a. Speichert die Daten dauerhaft
        System-->>Einlagerer: 6a. Aktualisiert die Ansicht
        System-->>Lagerhalter: 7a. Uebermittelt den Serviceauftrag an den Lagerhalter
    else Service abbuchen
        Einlagerer->>System: 3b. wählt einen gebuchten Service zur Abwahl aus
        Einlagerer->>System: 4b. Bestätigt die Auswahl
        note over System: Vertrag CO_UC_E6_2: removeBookedService
        System->>System: 5b. Entfernt die Servicebuchung dauerhaft
        System-->>Einlagerer: 6b. Aktualisiert die Ansicht
    end
\end{lstlisting}

\newpage

% ------------------------------------------------------------------------------
% UC_E7 – Fahrzeug registrieren
% ------------------------------------------------------------------------------
\phantomsection
\addcontentsline{toc}{subsubsection}{UC\_E7 – Fahrzeug registrieren}
\subsubsection*{UC\_E7 – Fahrzeug registrieren}

\begin{lstlisting}[basicstyle=\ttfamily\small, breaklines=true]
sequenceDiagram
    actor Einlagerer
    participant System

    Einlagerer->>System: 1. Ruft die Ansicht "Fahrzeugregistrierung" auf
    Einlagerer->>System: 2. Trägt die Fahrzeugdaten (Kennzeichen) ein
    Einlagerer->>System: 3. Bestätigt die Eingaben
    note over System: Vertrag CO_UC_E7_1: createCar
    System->>System: 4. Prueft Kennzeichen auf Gueltigkeit

    alt 4.1 Ausnahmeszenario: Kennzeichen bereits registriert
        System-->>Einlagerer: 4.1.1 Informiert den Einlagerer ueber den Konflikt
                                     und bittet um Korrektur
    end

    System->>System: 5. Initialisiert den Fahrzeugstatus (Fahrbereitschaft, Wartungsstand)
    System->>System: 6. Speichert die Daten dauerhaft
    System-->>Einlagerer: 7. Bestätigt dem Einlagerer die Registrierung
\end{lstlisting}

\newpage

% ------------------------------------------------------------------------------
% UC_E8 – Fahrzeuge verwalten
% ------------------------------------------------------------------------------
\phantomsection
\addcontentsline{toc}{subsubsection}{UC\_E8 – Fahrzeuge verwalten}
\subsubsection*{UC\_E8 – Fahrzeuge verwalten}

\begin{lstlisting}[basicstyle=\ttfamily\small, breaklines=true]
sequenceDiagram
    actor Einlagerer
    participant System

    Einlagerer->>System: 1. Ruft die Ansicht "Fahrzeuguebersicht" auf
    note over System: Vertrag CO_UC_E8_1: getAllCars
    System-->>Einlagerer: 2. Zeigt die Uebersicht der eigenen Fahrzeuge an

    alt Fahrzeug bearbeiten
        Einlagerer->>System: 3a. Passt die Fahrzeugdaten (Kennzeichen) an
        Einlagerer->>System: 4a. Bestätigt die Eingaben
        note over System: Vertrag CO_UC_E8_2: updateCar
        System->>System: 5a. Prueft Kennzeichen auf Gueltigkeit

        alt 5.1 Ausnahmeszenario: Kennzeichen bereits registriert
            System-->>Einlagerer: 5.1.1 Informiert den Einlagerer ueber den Konflikt
                                         und bittet um Korrektur
        end

        System->>System: 6a. Speichert die Daten dauerhaft
        System-->>Einlagerer: 7a. Aktualisiert die Ansicht
    else Fahrzeug loeschen
        Einlagerer->>System: 3b. wählt ein Fahrzeug zum Loeschen aus
        Einlagerer->>System: 4b. Bestätigt die Auswahl
        note over System: Vertrag CO_UC_E8_3: removeCar
        System->>System: 5b. Entfernt das Fahrzeug dauerhaft
        System-->>Einlagerer: 6b. Aktualisiert die Ansicht
    end
\end{lstlisting}

\newpage


% ==============================================================================
% TEIL 2: OPERATIONELLE VERTRäGE
% ==============================================================================

% ------------------------------------------------------------------------------
% CO_UC_B1_1 – createStorageOwner
% ------------------------------------------------------------------------------
\phantomsection
\addcontentsline{toc}{subsubsection}{CO\_UC\_B1\_1}
\begin{longtable}{|L|R|}
\hline
\rowcolor{tableHeader}
\textbf{Contract} & \textbf{CO\_UC\_B1\_1: createStorageOwner} \\ \hline
Operation & createStorageOwner(email, password, companyName, costs, serviceTime, location, storage, carManufacturers) \\ \hline
Querverweis & UC\_B1 – Lagerhalter registrieren \\ \hline
Nachbedingung &
\begin{itemize}[label=\checkmark, nosep, leftmargin=*]
    \item Eine Instanz \underline{so} von StorageOwner wurde erzeugt.
    \item \underline{so}.email wurde eine E-Mail-Adresse zugewiesen.
    \item \underline{so}.password wurde ein Passwort zugewiesen.
    \item \underline{so}.companyName wurde ein Firmenname zugewiesen.
    \item \underline{so}.costs wurde ein Kostenbetrag zugewiesen.
    \item \underline{so}.serviceTime wurde eine Servicezeit zugewiesen.
    \item Eine Instanz \underline{l} von Location wurde erzeugt und \underline{so}.location zugewiesen.
    \item Eine Instanz \underline{st} von Storage wurde erzeugt und \underline{so}.storage zugewiesen.
    \item \underline{so}.carManufacturers wurden Fahrzeughersteller zugewiesen.
\end{itemize} \\ \hline
\end{longtable}
\newpage

% ------------------------------------------------------------------------------
% CO_UC_B1_2 – createCarOwner
% ------------------------------------------------------------------------------
\phantomsection
\addcontentsline{toc}{subsubsection}{CO\_UC\_B1\_2}
\begin{longtable}{|L|R|}
\hline
\rowcolor{tableHeader}
\textbf{Contract} & \textbf{CO\_UC\_B1\_2: createCarOwner} \\ \hline
Operation & createCarOwner(email, password, firstName, lastName) \\ \hline
Querverweis & UC\_B1 – Einlagerer registrieren \\ \hline
Nachbedingung &
\begin{itemize}[label=\checkmark, nosep, leftmargin=*]
    \item Eine Instanz \underline{co} von CarOwner wurde erzeugt.
    \item \underline{co}.email wurde eine E-Mail-Adresse zugewiesen.
    \item \underline{co}.password wurde ein Passwort zugewiesen.
    \item \underline{co}.firstName wurde ein Vorname zugewiesen.
    \item \underline{co}.lastName wurde ein Nachname zugewiesen.
\end{itemize} \\ \hline
\end{longtable}
\newpage

% ------------------------------------------------------------------------------
% CO_UC_B2_1 – authenticate
% ------------------------------------------------------------------------------
\phantomsection
\addcontentsline{toc}{subsubsection}{CO\_UC\_B2\_1}
\begin{longtable}{|L|R|}
\hline
\rowcolor{tableHeader}
\textbf{Contract} & \textbf{CO\_UC\_B2\_1: authenticate} \\ \hline
Operation & authenticate(email, password) \\ \hline
Querverweis & UC\_B2 – Benutzer anmelden \\ \hline
Nachbedingung &
\begin{itemize}[label=\checkmark, nosep, leftmargin=*]
    \item Eine Instanz u von User wurde über email und password abgerufen und authentifiziert.
\end{itemize} \\ \hline
\end{longtable}
\newpage

% ------------------------------------------------------------------------------
% CO_UC_B3_1 (UC_B3.1) – getStorageOwner
% ------------------------------------------------------------------------------
\phantomsection
\addcontentsline{toc}{subsubsection}{CO\_UC\_B3\_1 (Lagerhalter verwalten)}
\begin{longtable}{|L|R|}
\hline
\rowcolor{tableHeader}
\textbf{Contract} & \textbf{CO\_UC\_B3\_1: getStorageOwner} \\ \hline
Operation & getStorageOwner(id) \\ \hline
Querverweis & UC\_B3.1 – Lagerhalter verwalten; UC\_L2 – Stellplatzkapazit{\"a}ten verwalten \\ \hline
Nachbedingung &
\begin{itemize}[label=\checkmark, nosep, leftmargin=*]
    \item Die aktuellen Profildaten von \underline{so} wurden abgerufen und zur{\"u}ckgegeben.
\end{itemize} \\ \hline
\end{longtable}
\newpage

% ------------------------------------------------------------------------------
% CO_UC_B3_1 (UC_B3.2) – getCarOwner
% ------------------------------------------------------------------------------
\phantomsection
\addcontentsline{toc}{subsubsection}{CO\_UC\_B3\_1 (Einlagerer verwalten)}
\begin{longtable}{|L|R|}
\hline
\rowcolor{tableHeader}
\textbf{Contract} & \textbf{CO\_UC\_B3\_1: getCarOwner} \\ \hline
Operation & getCarOwner(id) \\ \hline
Querverweis & UC\_B3.2 – Einlagerer verwalten \\ \hline
Nachbedingung &
\begin{itemize}[label=\checkmark, nosep, leftmargin=*]
    \item Die aktuellen Profildaten von \underline{co} wurden abgerufen und zur{\"u}ckgegeben.
\end{itemize} \\ \hline
\end{longtable}
\newpage

% ------------------------------------------------------------------------------
% CO_UC_B3_2 (UC_B3.1) – updateStorageOwner
% ------------------------------------------------------------------------------
\phantomsection
\addcontentsline{toc}{subsubsection}{CO\_UC\_B3\_2 (Lagerhalter verwalten)}
\begin{longtable}{|L|R|}
\hline
\rowcolor{tableHeader}
\textbf{Contract} & \textbf{CO\_UC\_B3\_2: updateStorageOwner} \\ \hline
Operation & updateStorageOwner(email, password, companyName, costs, serviceTime, location, storage, carManufacturers) \\ \hline
Querverweis & UC\_B3.1 – Lagerhalter verwalten; UC\_L2 – Stellplatzkapazit{\"a}ten verwalten \\ \hline
Nachbedingung &
\begin{itemize}[label=\checkmark, nosep, leftmargin=*]
    \item \underline{so}.email wurde aktualisiert.
    \item \underline{so}.password wurde aktualisiert.
    \item \underline{so}.companyName wurde aktualisiert.
    \item \underline{so}.costs wurde aktualisiert.
    \item \underline{so}.serviceTime wurde aktualisiert.
    \item \underline{so}.location wurde aktualisiert.
    \item \underline{so}.storage wurde aktualisiert.
    \item \underline{so}.carManufacturers wurde aktualisiert.
\end{itemize} \\ \hline
\end{longtable}
\newpage

% ------------------------------------------------------------------------------
% CO_UC_B3_2 (UC_B3.2) – updateCarOwner
% ------------------------------------------------------------------------------
\phantomsection
\addcontentsline{toc}{subsubsection}{CO\_UC\_B3\_2 (Einlagerer verwalten)}
\begin{longtable}{|L|R|}
\hline
\rowcolor{tableHeader}
\textbf{Contract} & \textbf{CO\_UC\_B3\_2: updateCarOwner} \\ \hline
Operation & updateCarOwner(email, password, firstName, lastName) \\ \hline
Querverweis & UC\_B3.2 – Einlagerer verwalten \\ \hline
Nachbedingung &
\begin{itemize}[label=\checkmark, nosep, leftmargin=*]
    \item \underline{co}.email wurde aktualisiert.
    \item \underline{co}.password wurde aktualisiert.
    \item \underline{co}.firstName wurde aktualisiert.
    \item \underline{co}.lastName wurde aktualisiert.
\end{itemize} \\ \hline
\end{longtable}
\newpage

% ------------------------------------------------------------------------------
% CO_UC_L1_1 – getAllServices
% ------------------------------------------------------------------------------
\phantomsection
\addcontentsline{toc}{subsubsection}{CO\_UC\_L1\_1}
\begin{longtable}{|L|R|}
\hline
\rowcolor{tableHeader}
\textbf{Contract} & \textbf{CO\_UC\_L1\_1: getAllServices} \\ \hline
Operation & getAllServices(storageOwnerId) \\ \hline
Querverweis & UC\_L1 – Serviceportfolio verwalten \\ \hline
Nachbedingung &
\begin{itemize}[label=\checkmark, nosep, leftmargin=*]
    \item Alle Services von \underline{so} wurden zur{\"u}ckgegeben.
\end{itemize} \\ \hline
\end{longtable}
\newpage

% ------------------------------------------------------------------------------
% CO_UC_L1_2 – createService
% ------------------------------------------------------------------------------
\phantomsection
\addcontentsline{toc}{subsubsection}{CO\_UC\_L1\_2}
\begin{longtable}{|L|R|}
\hline
\rowcolor{tableHeader}
\textbf{Contract} & \textbf{CO\_UC\_L1\_2: createService} \\ \hline
Operation & createService(title, description, price, isOptional) \\ \hline
Querverweis & UC\_L1 – Serviceportfolio verwalten; UC\_L6 – Optionale Services verwalten; UC\_E6 – Optionale Lagerservices buchen \\ \hline
Nachbedingung &
\begin{itemize}[label=\checkmark, nosep, leftmargin=*]
    \item Eine Instanz \underline{s} von Service wurde erzeugt.
    \item \underline{s}.title wurde ein Titel zugewiesen.
    \item \underline{s}.description wurde eine Beschreibung zugewiesen.
    \item \underline{s}.price wurde ein Preis zugewiesen.
    \item \underline{s}.isOptional wurde gesetzt.
    \item \underline{so}.services wurde um \underline{s} erweitert.
\end{itemize} \\ \hline
\end{longtable}
\newpage

% ------------------------------------------------------------------------------
% CO_UC_L1_3 – updateService
% ------------------------------------------------------------------------------
\phantomsection
\addcontentsline{toc}{subsubsection}{CO\_UC\_L1\_3}
\begin{longtable}{|L|R|}
\hline
\rowcolor{tableHeader}
\textbf{Contract} & \textbf{CO\_UC\_L1\_3: updateService} \\ \hline
Operation & updateService(title, description, price, isOptional) \\ \hline
Querverweis & UC\_L1 – Serviceportfolio verwalten; UC\_L6 – Optionale Services verwalten \\ \hline
Nachbedingung &
\begin{itemize}[label=\checkmark, nosep, leftmargin=*]
    \item \underline{s}.title wurde aktualisiert.
    \item \underline{s}.description wurde aktualisiert.
    \item \underline{s}.price wurde aktualisiert.
    \item \underline{s}.isOptional wurde aktualisiert.
\end{itemize} \\ \hline
\end{longtable}
\newpage

% ------------------------------------------------------------------------------
% CO_UC_L1_4 – deleteService
% ------------------------------------------------------------------------------
\phantomsection
\addcontentsline{toc}{subsubsection}{CO\_UC\_L1\_4}
\begin{longtable}{|L|R|}
\hline
\rowcolor{tableHeader}
\textbf{Contract} & \textbf{CO\_UC\_L1\_4: deleteService} \\ \hline
Operation & deleteService(serviceId) \\ \hline
Querverweis & UC\_L1 – Serviceportfolio verwalten; UC\_L6 – Optionale Services verwalten \\ \hline
Nachbedingung &
\begin{itemize}[label=\checkmark, nosep, leftmargin=*]
    \item \underline{s} wurde aus \underline{so}.services entfernt.
\end{itemize} \\ \hline
\end{longtable}
\newpage

% ------------------------------------------------------------------------------
% CO_UC_L3_1 – getStoredCars
% ------------------------------------------------------------------------------
\phantomsection
\addcontentsline{toc}{subsubsection}{CO\_UC\_L3\_1}
\begin{longtable}{|L|R|}
\hline
\rowcolor{tableHeader}
\textbf{Contract} & \textbf{CO\_UC\_L3\_1: getStoredCars} \\ \hline
Operation & getStoredCars(storageOwnerId) \\ \hline
Querverweis & UC\_L3 – Fahrzeugstatus aktualisieren \\ \hline
Nachbedingung &
\begin{itemize}[label=\checkmark, nosep, leftmargin=*]
    \item Alle eingelagerten Fahrzeuge von \underline{so} wurden zur{\"u}ckgegeben.
\end{itemize} \\ \hline
\end{longtable}
\newpage

% ------------------------------------------------------------------------------
% CO_UC_L3_2 – updateStoredCarStatus
% ------------------------------------------------------------------------------
\phantomsection
\addcontentsline{toc}{subsubsection}{CO\_UC\_L3\_2}
\begin{longtable}{|L|R|}
\hline
\rowcolor{tableHeader}
\textbf{Contract} & \textbf{CO\_UC\_L3\_2: updateStoredCarStatus} \\ \hline
Operation & updateStoredCarStatus(readyToDrive, isMaintained) \\ \hline
Querverweis & UC\_L3 – Fahrzeugstatus aktualisieren \\ \hline
Nachbedingung &
\begin{itemize}[label=\checkmark, nosep, leftmargin=*]
    \item \underline{c}.readyToDrive wurde auf den neuen Wert gesetzt.
    \item \underline{c}.isMaintained wurde auf den neuen Wert gesetzt.
\end{itemize} \\ \hline
\end{longtable}
\newpage

% ------------------------------------------------------------------------------
% CO_UC_L6_1 – getOptionalServices
% ------------------------------------------------------------------------------
\phantomsection
\addcontentsline{toc}{subsubsection}{CO\_UC\_L6\_1}
\begin{longtable}{|L|R|}
\hline
\rowcolor{tableHeader}
\textbf{Contract} & \textbf{CO\_UC\_L6\_1: getOptionalServices} \\ \hline
Operation & getOptionalServices(storageOwnerId) \\ \hline
Querverweis & UC\_L6 – Optionale Services verwalten; UC\_E6 – Optionale Lagerservices buchen \\ \hline
Nachbedingung &
\begin{itemize}[label=\checkmark, nosep, leftmargin=*]
    \item Alle optionalen Services von \underline{so} wurden zur{\"u}ckgegeben.
\end{itemize} \\ \hline
\end{longtable}
\newpage

% ------------------------------------------------------------------------------
% CO_UC_L7_1 – getIncomingRequests
% ------------------------------------------------------------------------------
\phantomsection
\addcontentsline{toc}{subsubsection}{CO\_UC\_L7\_1}
\begin{longtable}{|L|R|}
\hline
\rowcolor{tableHeader}
\textbf{Contract} & \textbf{CO\_UC\_L7\_1: getIncomingRequests} \\ \hline
Operation & getIncomingRequests(storageOwnerId) \\ \hline
Querverweis & UC\_L7 – Stellplatzanfragen verwalten \\ \hline
Nachbedingung &
\begin{itemize}[label=\checkmark, nosep, leftmargin=*]
    \item Alle eingehenden Stellplatzanfragen von \underline{so} wurden zur{\"u}ckgegeben.
\end{itemize} \\ \hline
\end{longtable}
\newpage

% ------------------------------------------------------------------------------
% CO_UC_L7_2 – createOffer
% ------------------------------------------------------------------------------
\phantomsection
\addcontentsline{toc}{subsubsection}{CO\_UC\_L7\_2}
\begin{longtable}{|L|R|}
\hline
\rowcolor{tableHeader}
\textbf{Contract} & \textbf{CO\_UC\_L7\_2: createOffer} \\ \hline
Operation & createOffer(requestId, description, costs) \\ \hline
Querverweis & UC\_L7 – Stellplatzanfragen verwalten \\ \hline
Nachbedingung &
\begin{itemize}[label=\checkmark, nosep, leftmargin=*]
    \item Eine Instanz \underline{o} von Offer wurde erzeugt.
    \item \underline{o}.description wurde eine Beschreibung zugewiesen.
    \item \underline{o}.costs wurde ein Kostenbetrag zugewiesen.
    \item \underline{o}.status wurde auf PENDING gesetzt.
    \item \underline{r}.offer wurde auf \underline{o} gesetzt.
    \item \underline{r}.status wurde auf CLOSED gesetzt.
    \item \underline{so}.offers wurde um \underline{o} erweitert.
\end{itemize} \\ \hline
\end{longtable}
\newpage

% ------------------------------------------------------------------------------
% CO_UC_L7_3 – getAllOffers
% ------------------------------------------------------------------------------
\phantomsection
\addcontentsline{toc}{subsubsection}{CO\_UC\_L7\_3}
\begin{longtable}{|L|R|}
\hline
\rowcolor{tableHeader}
\textbf{Contract} & \textbf{CO\_UC\_L7\_3: getAllOffers} \\ \hline
Operation & getAllOffers(storageOwnerId) \\ \hline
Querverweis & UC\_L7 – Stellplatzanfragen verwalten \\ \hline
Nachbedingung &
\begin{itemize}[label=\checkmark, nosep, leftmargin=*]
    \item Alle Angebote von \underline{so} wurden zur{\"u}ckgegeben.
\end{itemize} \\ \hline
\end{longtable}
\newpage

% ------------------------------------------------------------------------------
% CO_UC_L7_4 – deleteOffer
% ------------------------------------------------------------------------------
\phantomsection
\addcontentsline{toc}{subsubsection}{CO\_UC\_L7\_4}
\begin{longtable}{|L|R|}
\hline
\rowcolor{tableHeader}
\textbf{Contract} & \textbf{CO\_UC\_L7\_4: deleteOffer} \\ \hline
Operation & deleteOffer(offerId) \\ \hline
Querverweis & UC\_L7 – Stellplatzanfragen verwalten \\ \hline
Nachbedingung &
\begin{itemize}[label=\checkmark, nosep, leftmargin=*]
    \item \underline{o} wurde aus \underline{so}.offers entfernt.
\end{itemize} \\ \hline
\end{longtable}
\newpage

% ------------------------------------------------------------------------------
% CO_UC_E1_1 – searchByCity
% ------------------------------------------------------------------------------
\phantomsection
\addcontentsline{toc}{subsubsection}{CO\_UC\_E1\_1}
\begin{longtable}{|L|R|}
\hline
\rowcolor{tableHeader}
\textbf{Contract} & \textbf{CO\_UC\_E1\_1: searchByCity} \\ \hline
Operation & searchByCity(city) \\ \hline
Querverweis & UC\_E1 – Stellplatz suchen; UC\_E2 – Stellplatzangebot anfordern \\ \hline
Nachbedingung &
\begin{itemize}[label=\checkmark, nosep, leftmargin=*]
    \item Alle Lagerhalter mit Standort in der gesuchten Stadt wurden zur{\"u}ckgegeben.
\end{itemize} \\ \hline
\end{longtable}
\newpage

% ------------------------------------------------------------------------------
% CO_UC_E2_1 – createRequest
% ------------------------------------------------------------------------------
\phantomsection
\addcontentsline{toc}{subsubsection}{CO\_UC\_E2\_1}
\begin{longtable}{|L|R|}
\hline
\rowcolor{tableHeader}
\textbf{Contract} & \textbf{CO\_UC\_E2\_1: createRequest} \\ \hline
Operation & createRequest(storageOwnerId, description) \\ \hline
Querverweis & UC\_E2 – Stellplatzangebot anfordern \\ \hline
Nachbedingung &
\begin{itemize}[label=\checkmark, nosep, leftmargin=*]
    \item Eine Instanz \underline{r} von Request wurde erzeugt.
    \item \underline{r}.type wurde auf STELLPLATZ gesetzt.
    \item \underline{r}.description wurde eine Beschreibung zugewiesen.
    \item \underline{r}.status wurde auf OPEN gesetzt.
    \item \underline{co}.requests wurde um \underline{r} erweitert.
    \item \underline{so}.incomingRequests wurde um \underline{r} erweitert.
\end{itemize} \\ \hline
\end{longtable}
\newpage

% ------------------------------------------------------------------------------
% CO_UC_E2_2 – getAllRequests
% ------------------------------------------------------------------------------
\phantomsection
\addcontentsline{toc}{subsubsection}{CO\_UC\_E2\_2}
\begin{longtable}{|L|R|}
\hline
\rowcolor{tableHeader}
\textbf{Contract} & \textbf{CO\_UC\_E2\_2: getAllRequests} \\ \hline
Operation & getAllRequests(carOwnerId) \\ \hline
Querverweis & UC\_E2 – Stellplatzangebot anfordern \\ \hline
Nachbedingung &
\begin{itemize}[label=\checkmark, nosep, leftmargin=*]
    \item Alle Stellplatzanfragen von \underline{co} wurden zur{\"u}ckgegeben.
\end{itemize} \\ \hline
\end{longtable}
\newpage

% ------------------------------------------------------------------------------
% CO_UC_E2_3 – deleteRequest
% ------------------------------------------------------------------------------
\phantomsection
\addcontentsline{toc}{subsubsection}{CO\_UC\_E2\_3}
\begin{longtable}{|L|R|}
\hline
\rowcolor{tableHeader}
\textbf{Contract} & \textbf{CO\_UC\_E2\_3: deleteRequest} \\ \hline
Operation & deleteRequest(requestId) \\ \hline
Querverweis & UC\_E2 – Stellplatzangebot anfordern \\ \hline
Nachbedingung &
\begin{itemize}[label=\checkmark, nosep, leftmargin=*]
    \item \underline{r} wurde aus \underline{co}.requests entfernt.
\end{itemize} \\ \hline
\end{longtable}
\newpage

% ------------------------------------------------------------------------------
% CO_UC_E3_1 – getReceivedOffers
% ------------------------------------------------------------------------------
\phantomsection
\addcontentsline{toc}{subsubsection}{CO\_UC\_E3\_1}
\begin{longtable}{|L|R|}
\hline
\rowcolor{tableHeader}
\textbf{Contract} & \textbf{CO\_UC\_E3\_1: getReceivedOffers} \\ \hline
Operation & getReceivedOffers(carOwnerId) \\ \hline
Querverweis & UC\_E3 – Stellplatzangebote verwalten \\ \hline
Nachbedingung &
\begin{itemize}[label=\checkmark, nosep, leftmargin=*]
    \item Alle empfangenen Angebote von \underline{co} wurden zur{\"u}ckgegeben.
\end{itemize} \\ \hline
\end{longtable}
\newpage

% ------------------------------------------------------------------------------
% CO_UC_E3_2 – respondToOffer
% ------------------------------------------------------------------------------
\phantomsection
\addcontentsline{toc}{subsubsection}{CO\_UC\_E3\_2}
\begin{longtable}{|L|R|}
\hline
\rowcolor{tableHeader}
\textbf{Contract} & \textbf{CO\_UC\_E3\_2: respondToOffer} \\ \hline
Operation & respondToOffer(offerId, accepted) \\ \hline
Querverweis & UC\_E3 – Stellplatzangebote verwalten \\ \hline
Nachbedingung &
\begin{itemize}[label=\checkmark, nosep, leftmargin=*]
    \item {[accepted = true]} \underline{o}.status wurde auf ACCEPTED gesetzt.
    \item {[accepted = false]} \underline{o}.status wurde auf REJECTED gesetzt.
\end{itemize} \\ \hline
\end{longtable}
\newpage

% ------------------------------------------------------------------------------
% CO_UC_E5_1 – createAppointment
% ------------------------------------------------------------------------------
\phantomsection
\addcontentsline{toc}{subsubsection}{CO\_UC\_E5\_1}
\begin{longtable}{|L|R|}
\hline
\rowcolor{tableHeader}
\textbf{Contract} & \textbf{CO\_UC\_E5\_1: createAppointment} \\ \hline
Operation & createAppointment(type, date, time) \\ \hline
Querverweis & UC\_E5 – Ein- \& Auslagerungstermin koordinieren \\ \hline
Nachbedingung &
\begin{itemize}[label=\checkmark, nosep, leftmargin=*]
    \item Eine Instanz \underline{a} von Appointment wurde erzeugt.
    \item \underline{a}.type wurde zugewiesen.
    \item \underline{a}.date wurde zugewiesen.
    \item \underline{a}.time wurde zugewiesen.
    \item \underline{a}.status wurde auf AVAILABLE gesetzt.
    \item \underline{so}.appointments wurde um \underline{a} erweitert.
\end{itemize} \\ \hline
\end{longtable}
\newpage

% ------------------------------------------------------------------------------
% CO_UC_E5_2 – getAvailableAppointments
% ------------------------------------------------------------------------------
\phantomsection
\addcontentsline{toc}{subsubsection}{CO\_UC\_E5\_2}
\begin{longtable}{|L|R|}
\hline
\rowcolor{tableHeader}
\textbf{Contract} & \textbf{CO\_UC\_E5\_2: getAvailableAppointments} \\ \hline
Operation & getAvailableAppointments(carOwnerId, storageOwnerId) \\ \hline
Querverweis & UC\_E5 – Ein- \& Auslagerungstermin koordinieren \\ \hline
Nachbedingung &
\begin{itemize}[label=\checkmark, nosep, leftmargin=*]
    \item Alle verf{\"u}gbaren Termine von \underline{so} wurden zur{\"u}ckgegeben.
\end{itemize} \\ \hline
\end{longtable}
\newpage

% ------------------------------------------------------------------------------
% CO_UC_E5_3 – bookAppointment
% ------------------------------------------------------------------------------
\phantomsection
\addcontentsline{toc}{subsubsection}{CO\_UC\_E5\_3}
\begin{longtable}{|L|R|}
\hline
\rowcolor{tableHeader}
\textbf{Contract} & \textbf{CO\_UC\_E5\_3: bookAppointment} \\ \hline
Operation & bookAppointment(storageOwnerId, carId, type, date, time) \\ \hline
Querverweis & UC\_E5 – Ein- \& Auslagerungstermin koordinieren \\ \hline
Nachbedingung &
\begin{itemize}[label=\checkmark, nosep, leftmargin=*]
    \item \underline{a}.status wurde auf BOOKED gesetzt.
    \item \underline{a}.carOwner wurde mit \underline{co} verkn{\"u}pft.
    \item \underline{a}.subject wurde mit \underline{c} verkn{\"u}pft.
    \item \underline{co}.bookedAppointments wurde um \underline{a} erweitert.
    \item {[type = EINLAGERUNG]} \underline{c}.currentStorage wurde auf \underline{so} gesetzt.
    \item {[type = EINLAGERUNG]} \underline{so}.storedCars wurde um \underline{c} erweitert.
    \item {[type = AUSLAGERUNG]} \underline{c}.currentStorage wurde auf null gesetzt.
    \item {[type = AUSLAGERUNG]} \underline{c} wurde aus \underline{so}.storedCars entfernt.
\end{itemize} \\ \hline
\end{longtable}
\newpage

% ------------------------------------------------------------------------------
% CO_UC_E5_4 – getAllBookedAppointments
% ------------------------------------------------------------------------------
\phantomsection
\addcontentsline{toc}{subsubsection}{CO\_UC\_E5\_4}
\begin{longtable}{|L|R|}
\hline
\rowcolor{tableHeader}
\textbf{Contract} & \textbf{CO\_UC\_E5\_4: getAllBookedAppointments} \\ \hline
Operation & getAllBookedAppointments(carOwnerId) \\ \hline
Querverweis & UC\_E5 – Ein- \& Auslagerungstermin koordinieren \\ \hline
Nachbedingung &
\begin{itemize}[label=\checkmark, nosep, leftmargin=*]
    \item Alle gebuchten Termine von \underline{co} wurden zur{\"u}ckgegeben.
\end{itemize} \\ \hline
\end{longtable}
\newpage

% ------------------------------------------------------------------------------
% CO_UC_E5_5 – cancelAppointment
% ------------------------------------------------------------------------------
\phantomsection
\addcontentsline{toc}{subsubsection}{CO\_UC\_E5\_5}
\begin{longtable}{|L|R|}
\hline
\rowcolor{tableHeader}
\textbf{Contract} & \textbf{CO\_UC\_E5\_5: cancelAppointment} \\ \hline
Operation & cancelAppointment(appointmentId) \\ \hline
Querverweis & UC\_E5 – Ein- \& Auslagerungstermin koordinieren \\ \hline
Nachbedingung &
\begin{itemize}[label=\checkmark, nosep, leftmargin=*]
    \item \underline{a}.status wurde auf CANCELLED gesetzt.
\end{itemize} \\ \hline
\end{longtable}
\newpage

% ------------------------------------------------------------------------------
% CO_UC_E5_6 – getAllAppointments
% ------------------------------------------------------------------------------
\phantomsection
\addcontentsline{toc}{subsubsection}{CO\_UC\_E5\_6}
\begin{longtable}{|L|R|}
\hline
\rowcolor{tableHeader}
\textbf{Contract} & \textbf{CO\_UC\_E5\_6: getAllAppointments} \\ \hline
Operation & getAllAppointments(storageOwnerId) \\ \hline
Querverweis & UC\_E5 – Ein- \& Auslagerungstermin koordinieren \\ \hline
Nachbedingung &
\begin{itemize}[label=\checkmark, nosep, leftmargin=*]
    \item Alle Termine von \underline{so} wurden zur{\"u}ckgegeben.
\end{itemize} \\ \hline
\end{longtable}
\newpage

% ------------------------------------------------------------------------------
% CO_UC_E5_7 – deleteAppointment
% ------------------------------------------------------------------------------
\phantomsection
\addcontentsline{toc}{subsubsection}{CO\_UC\_E5\_7}
\begin{longtable}{|L|R|}
\hline
\rowcolor{tableHeader}
\textbf{Contract} & \textbf{CO\_UC\_E5\_7: deleteAppointment} \\ \hline
Operation & deleteAppointment(appointmentId) \\ \hline
Querverweis & UC\_E5 – Ein- \& Auslagerungstermin koordinieren \\ \hline
Nachbedingung &
\begin{itemize}[label=\checkmark, nosep, leftmargin=*]
    \item \underline{a} wurde aus \underline{so}.appointments entfernt.
\end{itemize} \\ \hline
\end{longtable}
\newpage

% ------------------------------------------------------------------------------
% CO_UC_E6_1 – bookOptionalService
% ------------------------------------------------------------------------------
\phantomsection
\addcontentsline{toc}{subsubsection}{CO\_UC\_E6\_1}
\begin{longtable}{|L|R|}
\hline
\rowcolor{tableHeader}
\textbf{Contract} & \textbf{CO\_UC\_E6\_1: bookOptionalService} \\ \hline
Operation & bookOptionalService(carId, storageOwnerId, serviceId) \\ \hline
Querverweis & UC\_E6 – Optionale Lagerservices buchen \\ \hline
Nachbedingung &
\begin{itemize}[label=\checkmark, nosep, leftmargin=*]
    \item \underline{s} wurde zu \underline{c}.bookedServices hinzugef{\"u}gt.
\end{itemize} \\ \hline
\end{longtable}
\newpage

% ------------------------------------------------------------------------------
% CO_UC_E6_2 – removeBookedService
% ------------------------------------------------------------------------------
\phantomsection
\addcontentsline{toc}{subsubsection}{CO\_UC\_E6\_2}
\begin{longtable}{|L|R|}
\hline
\rowcolor{tableHeader}
\textbf{Contract} & \textbf{CO\_UC\_E6\_2: removeBookedService} \\ \hline
Operation & removeBookedService(carId, serviceId) \\ \hline
Querverweis & UC\_E6 – Optionale Lagerservices buchen \\ \hline
Nachbedingung &
\begin{itemize}[label=\checkmark, nosep, leftmargin=*]
    \item \underline{s} wurde aus \underline{c}.bookedServices entfernt.
\end{itemize} \\ \hline
\end{longtable}
\newpage

% ------------------------------------------------------------------------------
% CO_UC_E7_1 – createCar
% ------------------------------------------------------------------------------
\phantomsection
\addcontentsline{toc}{subsubsection}{CO\_UC\_E7\_1}
\begin{longtable}{|L|R|}
\hline
\rowcolor{tableHeader}
\textbf{Contract} & \textbf{CO\_UC\_E7\_1: createCar} \\ \hline
Operation & createCar(licensePlate) \\ \hline
Querverweis & UC\_E7 – Fahrzeug registrieren \\ \hline
Nachbedingung &
\begin{itemize}[label=\checkmark, nosep, leftmargin=*]
    \item Eine Instanz \underline{c} von Car wurde erzeugt.
    \item \underline{c}.licensePlate wurde ein Kennzeichen zugewiesen.
    \item \underline{c}.readyToDrive wurde auf true gesetzt.
    \item \underline{c}.isMaintained wurde auf false gesetzt.
    \item \underline{co}.cars wurde um \underline{c} erweitert.
\end{itemize} \\ \hline
\end{longtable}
\newpage

% ------------------------------------------------------------------------------
% CO_UC_E8_1 – getAllCars
% ------------------------------------------------------------------------------
\phantomsection
\addcontentsline{toc}{subsubsection}{CO\_UC\_E8\_1}
\begin{longtable}{|L|R|}
\hline
\rowcolor{tableHeader}
\textbf{Contract} & \textbf{CO\_UC\_E8\_1: getAllCars} \\ \hline
Operation & getAllCars(carOwnerId) \\ \hline
Querverweis & UC\_E8 – Fahrzeuge verwalten \\ \hline
Nachbedingung &
\begin{itemize}[label=\checkmark, nosep, leftmargin=*]
    \item Alle Fahrzeuge von \underline{co} wurden zur{\"u}ckgegeben.
\end{itemize} \\ \hline
\end{longtable}
\newpage

% ------------------------------------------------------------------------------
% CO_UC_E8_2 – updateCar
% ------------------------------------------------------------------------------
\phantomsection
\addcontentsline{toc}{subsubsection}{CO\_UC\_E8\_2}
\begin{longtable}{|L|R|}
\hline
\rowcolor{tableHeader}
\textbf{Contract} & \textbf{CO\_UC\_E8\_2: updateCar} \\ \hline
Operation & updateCar(licensePlate) \\ \hline
Querverweis & UC\_E8 – Fahrzeuge verwalten \\ \hline
Nachbedingung &
\begin{itemize}[label=\checkmark, nosep, leftmargin=*]
    \item \underline{c}.licensePlate wurde auf das neue Kennzeichen aktualisiert.
\end{itemize} \\ \hline
\end{longtable}
\newpage

% ------------------------------------------------------------------------------
% CO_UC_E8_3 – removeCar
% ------------------------------------------------------------------------------
\phantomsection
\addcontentsline{toc}{subsubsection}{CO\_UC\_E8\_3}
\begin{longtable}{|L|R|}
\hline
\rowcolor{tableHeader}
\textbf{Contract} & \textbf{CO\_UC\_E8\_3: removeCar} \\ \hline
Operation & removeCar(carId) \\ \hline
Querverweis & UC\_E8 – Fahrzeuge verwalten \\ \hline
Nachbedingung &
\begin{itemize}[label=\checkmark, nosep, leftmargin=*]
    \item \underline{c} wurde aus \underline{co}.cars entfernt.
\end{itemize} \\ \hline
\end{longtable}
\newpage


% ==============================================================================
% TEIL 3: ENTWURFSKLASSENDIAGRAMM
% ==============================================================================
\phantomsection
\addcontentsline{toc}{subsection}{Entwurfsklassendiagramm}
\subsection*{Entwurfsklassendiagramm}

\begin{lstlisting}[basicstyle=\ttfamily\small, breaklines=true]
---
config:
  layout: elk
---
classDiagram
    class CarOwnerController {
        +createCarOwner(email, password, firstName, lastName)
        +getCarOwner(id) CarOwner
        +updateCarOwner(email, password, firstName, lastName)
        +searchByCity(city) List~StorageOwner~
        +createCar(licensePlate) Car
        +getAllCars(carOwner) List~Car~
        +updateCar(carId, licensePlate) Car
        +removeCar(carId)
        +bookOptionalService(carId, storageOwnerId, serviceId) Car
        +removeBookedService(carId, serviceId)
        +createRequest(storageOwnerId, description) Request
        +getAllRequests(carOwner) List~Request~
        +deleteRequest(requestId)
        +getReceivedOffers(carOwner) List~Offer~
        +respondToOffer(offerId, accepted) Offer
        +getAvailableAppointments(carOwnerId, storageOwnerId) List~Appointment~
        +bookAppointment(storageOwnerId, carId, type, date, time) Appointment
        +getAllBookedAppointments(carOwner) List~Appointment~
        +cancelAppointment(appointmentId)
    }
    class StorageOwnerController {
        +createStorageOwner(email, password, companyName, costs, serviceTime, location, storage, carManufacturers)
        +getStorageOwner(id) StorageOwner
        +updateStorageOwner(email, password, companyName, costs, serviceTime, location, storage, carManufacturers)
        +getAllServices(storageOwner) List~Service~
        +createService(title, description, price, isOptional) Service
        +updateService(serviceId, title, description, price, isOptional) Service
        +deleteService(serviceId)
        +getOptionalServices(storageOwner) List~Service~
        +getStoredCars(storageOwner) List~Car~
        +updateStoredCarStatus(carId, readyToDrive, isMaintained) Car
        +getIncomingRequests(storageOwner) List~Request~
        +createOffer(requestId, description, costs) Offer
        +getAllOffers(storageOwner) List~Offer~
        +deleteOffer(offerId)
        +createAppointment(type, date, time) Appointment
        +getAllAppointments(storageOwner) List~Appointment~
        +deleteAppointment(appointmentId)
    }
    class AuthController{
        +authenticate(email, password)
    }
    class User {
        #email : String
        #password : String
        +changeEmail(email) void
        +changePassword(password) void
    }
    class CarOwner {
        +firstName : String
        +lastName : String
        +updateProfile(email, firstName, lastName) void
        +registerCar(licensePlate) Car
        +submitRequest(storageOwner, description) Request
        +removeCar(car) void
        +removeRequest(request) void
        +addAppointment(appointment) void
    }
    class StorageOwner {
        +companyName : String
        +costs : float
        +serviceTime : LocalTime
        +carManufacturers : Set~String~
        +updateProfile(email, companyName, costs, serviceTime, location, storage, carManufacturers) void
        +openService(title, description, price, isOptional) Service
        +makeOffer(request, description, costs) Offer
        +scheduleAppointment(type, date, time) Appointment
        +removeService(service) void
        +addCarToStorage(car) void
        +removeCarFromStorage(car) void
        +addRequest(request) void
        +removeOffer(offer) void
        +removeAppointment(appointment) void
    }
    class Location {
        +street : String
        +houseNumber : String
        +postalCode : String
        +city : String
    }
    class Storage {
        +conditions : List~String~
        +capacity : int
    }
    class Car {
        +licensePlate : String
        +readyToDrive : boolean
        +isMaintained : boolean
        +updateStatus(readyToDrive, isMaintained) void
        +changeLicensePlate(licensePlate) void
        +relocateTo(storageOwner) void
        +addBookedOptionalService(service) void
        +removeBookedOptionalService(service) void
        ~setCurrentStorage(storageOwner) void
    }
    class Service {
        +title : String
        +description : String
        +price : float
        +isOptional : boolean
        +updateDetails(title, description, price, isOptional) void
    }
    class Request {
        +type : RequestType
        +description : String
        +status : RequestStatus
        +addOffer(offer) void
    }
    class Offer {
        +description : String
        +costs : float
        +status : OfferStatus
        +accept() void
        +reject() void
    }
    class Appointment {
        +type : AppointmentType
        +status : AppointmentStatus
        +date : LocalDateTime
        +time : LocalTime
        +book(carOwner, car) void
        +cancel() void
    }
    class AppointmentType {EINLAGERUNG
    AUSLAGERUNG }
    class AppointmentStatus {AVAILABLE
    BOOKED
    CANCELLED }
    class OfferStatus {PENDING
    ACCEPTED
    REJECTED }
    class RequestStatus {OPEN
    CLOSED }
    class RequestType {STELLPLATZ }


    <<enumeration>> AppointmentType
    <<enumeration>> AppointmentStatus
    <<enumeration>> OfferStatus
    <<enumeration>> RequestStatus
    <<enumeration>> RequestType

    User <|-- CarOwner
    User <|-- StorageOwner
    CarOwner "1" --> "*" Car : cars
    CarOwner "1" --> "*" Request : requests
    CarOwner "1" --> "*" Appointment : bookedAppointments
    StorageOwner "1" --> "*" Service : services
    StorageOwner "1" --> "*" Offer : offers
    StorageOwner "1" --> "*" Appointment : appointments
    StorageOwner "1" --> "*" Request : incomingRequests
    Car "*" --> "*" Service : bookedServices
    Car "*" --> "0..1" StorageOwner : currentStorage
    Request "1" --> "0..1" Offer : offer
    Appointment "*" --> "0..1" Car : subject
    StorageOwner "1" *-- "1" Location : location
    StorageOwner "1" *-- "1" Storage : storage
    Request --> RequestType
    Request --> RequestStatus
    Offer --> OfferStatus
    Appointment --> AppointmentType
    Appointment --> AppointmentStatus
    CarOwnerController ..> CarOwner : delegiert
    StorageOwnerController ..> StorageOwner : delegiert
    AuthController..>User : delegiert
\end{lstlisting}

\newpage


% ==============================================================================
% TEIL 4: UML-ENTWURFSSEQUENZDIAGRAMME
% ==============================================================================

% ------------------------------------------------------------------------------
% ESD CO_UC_B1_1 – createStorageOwner
% ------------------------------------------------------------------------------
\phantomsection
\addcontentsline{toc}{subsubsection}{ESD CO\_UC\_B1\_1 – createStorageOwner}
\subsubsection*{ESD CO\_UC\_B1\_1 – createStorageOwner}

\begin{lstlisting}[basicstyle=\ttfamily\small, breaklines=true]
sequenceDiagram
    autonumber
    actor Lagerhalter
    participant ctrl as :StorageOwnerController
    participant so as so:StorageOwner
    participant l as l:Location
    participant st as st:Storage

    Lagerhalter->>ctrl: createStorageOwner(email, password, companyName,
                        costs, serviceTime, location, storage, carManufacturers)
    ctrl->>l: Location(street, houseNumber, postalCode, city)
    ctrl->>st: Storage(conditions, capacity)
    ctrl->>so: StorageOwner(email, password, companyName, costs, serviceTime,
                            l, st, carManufacturers)
\end{lstlisting}

\newpage

% ------------------------------------------------------------------------------
% ESD CO_UC_B1_2 – createCarOwner
% ------------------------------------------------------------------------------
\phantomsection
\addcontentsline{toc}{subsubsection}{ESD CO\_UC\_B1\_2 – createCarOwner}
\subsubsection*{ESD CO\_UC\_B1\_2 – createCarOwner}

\begin{lstlisting}[basicstyle=\ttfamily\small, breaklines=true]
sequenceDiagram
    autonumber
    actor Einlagerer
    participant ctrl as :CarOwnerController
    participant co as co:CarOwner

    Einlagerer->>ctrl: createCarOwner(email, password, firstName, lastName)
    ctrl->>co: CarOwner(email, password, firstName, lastName)
\end{lstlisting}

\newpage

% ------------------------------------------------------------------------------
% ESD CO_UC_B2_1 – authenticate
% ------------------------------------------------------------------------------
\phantomsection
\addcontentsline{toc}{subsubsection}{ESD CO\_UC\_B2\_1 – authenticate}
\subsubsection*{ESD CO\_UC\_B2\_1 – authenticate}

\begin{lstlisting}[basicstyle=\ttfamily\small, breaklines=true]
sequenceDiagram
    autonumber
    actor Benutzer
    participant a as :AuthController

    Benutzer->>a: authenticate(email, password)
    a-->>Benutzer: Anmeldung bestätigt
\end{lstlisting}

\newpage

% ------------------------------------------------------------------------------
% ESD CO_UC_B3_1 (Lagerhalter) – getStorageOwner
% ------------------------------------------------------------------------------
\phantomsection
\addcontentsline{toc}{subsubsection}{ESD CO\_UC\_B3\_1 – getStorageOwner}
\subsubsection*{ESD CO\_UC\_B3\_1 – getStorageOwner (UC\_B3.1 / UC\_L2)}

\begin{lstlisting}[basicstyle=\ttfamily\small, breaklines=true]
sequenceDiagram
    autonumber
    actor Lagerhalter
    participant ctrl as :StorageOwnerController

    Lagerhalter->>ctrl: getStorageOwner(id)
    ctrl-->>Lagerhalter: so:StorageOwner
\end{lstlisting}

\newpage

% ------------------------------------------------------------------------------
% ESD CO_UC_B3_1 (Einlagerer) – getCarOwner
% ------------------------------------------------------------------------------
\phantomsection
\addcontentsline{toc}{subsubsection}{ESD CO\_UC\_B3\_1 – getCarOwner}
\subsubsection*{ESD CO\_UC\_B3\_1 – getCarOwner (UC\_B3.2)}

\begin{lstlisting}[basicstyle=\ttfamily\small, breaklines=true]
sequenceDiagram
    autonumber
    actor Einlagerer
    participant ctrl as :CarOwnerController

    Einlagerer->>ctrl: getCarOwner(id)
    ctrl-->>Einlagerer: co:CarOwner
\end{lstlisting}

\newpage

% ------------------------------------------------------------------------------
% ESD CO_UC_B3_2 (Lagerhalter) – updateStorageOwner
% ------------------------------------------------------------------------------
\phantomsection
\addcontentsline{toc}{subsubsection}{ESD CO\_UC\_B3\_2 – updateStorageOwner}
\subsubsection*{ESD CO\_UC\_B3\_2 – updateStorageOwner (UC\_B3.1 / UC\_L2)}

\begin{lstlisting}[basicstyle=\ttfamily\small, breaklines=true]
sequenceDiagram
    autonumber
    actor Lagerhalter
    participant ctrl as :StorageOwnerController
    participant so as so:StorageOwner
    participant l as l:Location
    participant st as st:Storage

    Lagerhalter->>ctrl: updateStorageOwner(email, password, companyName,
                        costs, serviceTime, location, storage, carManufacturers)
    ctrl->>l: Location(street, houseNumber, postalCode, city)
    ctrl->>st: Storage(conditions, capacity)
    ctrl->>so: so.updateProfile(email, companyName, costs, serviceTime,
                        l, st, carManufacturers)
    ctrl->>so: so.changePassword(password)
\end{lstlisting}

\newpage

% ------------------------------------------------------------------------------
% ESD CO_UC_B3_2 (Einlagerer) – updateCarOwner
% ------------------------------------------------------------------------------
\phantomsection
\addcontentsline{toc}{subsubsection}{ESD CO\_UC\_B3\_2 – updateCarOwner}
\subsubsection*{ESD CO\_UC\_B3\_2 – updateCarOwner (UC\_B3.2)}

\begin{lstlisting}[basicstyle=\ttfamily\small, breaklines=true]
sequenceDiagram
    autonumber
    actor Einlagerer
    participant ctrl as :CarOwnerController
    participant co as co:CarOwner

    Einlagerer->>ctrl: updateCarOwner(email, password, firstName, lastName)
    ctrl->>co: co.updateProfile(email, firstName, lastName)
    ctrl->>co: co.changePassword(password)
\end{lstlisting}

\newpage

% ------------------------------------------------------------------------------
% ESD CO_UC_L1_1 – getAllServices
% ------------------------------------------------------------------------------
\phantomsection
\addcontentsline{toc}{subsubsection}{ESD CO\_UC\_L1\_1 – getAllServices}
\subsubsection*{ESD CO\_UC\_L1\_1 – getAllServices}

\begin{lstlisting}[basicstyle=\ttfamily\small, breaklines=true]
sequenceDiagram
    autonumber
    actor Lagerhalter
    participant ctrl as :StorageOwnerController
    participant so as so:StorageOwner

    Lagerhalter->>ctrl: getAllServices(storageOwner)
    ctrl->>so: so.getServices()
    so-->>ctrl: services
    ctrl-->>Lagerhalter: Liste von Services
\end{lstlisting}

\newpage

% ------------------------------------------------------------------------------
% ESD CO_UC_L1_2 – createService
% ------------------------------------------------------------------------------
\phantomsection
\addcontentsline{toc}{subsubsection}{ESD CO\_UC\_L1\_2 – createService}
\subsubsection*{ESD CO\_UC\_L1\_2 – createService}

\begin{lstlisting}[basicstyle=\ttfamily\small, breaklines=true]
sequenceDiagram
    autonumber
    actor Lagerhalter
    participant ctrl as :StorageOwnerController
    participant so as so:StorageOwner
    participant s as s:Service

    Lagerhalter->>ctrl: createService(title, description, price, isOptional)
    ctrl->>so: so.openService(title, description, price, isOptional)
    so->>s: Service(so, title, description, price, isOptional)
    so->>so: services.add(s)
\end{lstlisting}

\newpage

% ------------------------------------------------------------------------------
% ESD CO_UC_L1_3 – updateService
% ------------------------------------------------------------------------------
\phantomsection
\addcontentsline{toc}{subsubsection}{ESD CO\_UC\_L1\_3 – updateService}
\subsubsection*{ESD CO\_UC\_L1\_3 – updateService}

\begin{lstlisting}[basicstyle=\ttfamily\small, breaklines=true]
sequenceDiagram
    autonumber
    actor Lagerhalter
    participant ctrl as :StorageOwnerController
    participant so as so:StorageOwner
    participant s as s:Service

    Lagerhalter->>ctrl: updateService(serviceId, title, description, price, isOptional)
    ctrl->>so: so.getServices()
    so-->>ctrl: services
    note over ctrl: wählt s mit id = serviceId
    ctrl->>s: s.updateDetails(title, description, price, isOptional)
\end{lstlisting}

\newpage

% ------------------------------------------------------------------------------
% ESD CO_UC_L1_4 – deleteService
% ------------------------------------------------------------------------------
\phantomsection
\addcontentsline{toc}{subsubsection}{ESD CO\_UC\_L1\_4 – deleteService}
\subsubsection*{ESD CO\_UC\_L1\_4 – deleteService}

\begin{lstlisting}[basicstyle=\ttfamily\small, breaklines=true]
sequenceDiagram
    autonumber
    actor Lagerhalter
    participant ctrl as :StorageOwnerController
    participant so as so:StorageOwner

    Lagerhalter->>ctrl: deleteService(serviceId)
    ctrl->>so: so.getServices()
    so-->>ctrl: services
    note over ctrl: wählt s mit id = serviceId
    ctrl->>so: so.removeService(s)
\end{lstlisting}

\newpage

% ------------------------------------------------------------------------------
% ESD CO_UC_L3_1 – getStoredCars
% ------------------------------------------------------------------------------
\phantomsection
\addcontentsline{toc}{subsubsection}{ESD CO\_UC\_L3\_1 – getStoredCars}
\subsubsection*{ESD CO\_UC\_L3\_1 – getStoredCars}

\begin{lstlisting}[basicstyle=\ttfamily\small, breaklines=true]
sequenceDiagram
    autonumber
    actor Lagerhalter
    participant ctrl as :StorageOwnerController
    participant so as so:StorageOwner

    Lagerhalter->>ctrl: getStoredCars(storageOwner)
    ctrl->>so: so.getStoredCars()
    so-->>ctrl: storedCars
    ctrl-->>Lagerhalter: Liste von Fahrzeugen
\end{lstlisting}

\newpage

% ------------------------------------------------------------------------------
% ESD CO_UC_L3_2 – updateStoredCarStatus
% ------------------------------------------------------------------------------
\phantomsection
\addcontentsline{toc}{subsubsection}{ESD CO\_UC\_L3\_2 – updateStoredCarStatus}
\subsubsection*{ESD CO\_UC\_L3\_2 – updateStoredCarStatus}

\begin{lstlisting}[basicstyle=\ttfamily\small, breaklines=true]
sequenceDiagram
    autonumber
    actor Lagerhalter
    participant ctrl as :StorageOwnerController
    participant so as so:StorageOwner
    participant c as c:Car

    Lagerhalter->>ctrl: updateStoredCarStatus(carId, readyToDrive, isMaintained)
    ctrl->>so: so.getStoredCars()
    so-->>ctrl: storedCars
    note over ctrl: wählt c mit id = carId
    ctrl->>c: c.updateStatus(readyToDrive, isMaintained)
\end{lstlisting}

\newpage

% ------------------------------------------------------------------------------
% ESD CO_UC_L6_1 – getOptionalServices
% ------------------------------------------------------------------------------
\phantomsection
\addcontentsline{toc}{subsubsection}{ESD CO\_UC\_L6\_1 – getOptionalServices}
\subsubsection*{ESD CO\_UC\_L6\_1 – getOptionalServices}

\begin{lstlisting}[basicstyle=\ttfamily\small, breaklines=true]
sequenceDiagram
    autonumber
    actor Lagerhalter
    participant ctrl as :StorageOwnerController
    participant so as so:StorageOwner

    Lagerhalter->>ctrl: getOptionalServices(storageOwner)
    ctrl->>so: so.getServices()
    so-->>ctrl: services
    note over ctrl: filtert Services mit isOptional = true
    ctrl-->>Lagerhalter: Liste von optionalen Services
\end{lstlisting}

\newpage

% ------------------------------------------------------------------------------
% ESD CO_UC_L7_1 – getIncomingRequests
% ------------------------------------------------------------------------------
\phantomsection
\addcontentsline{toc}{subsubsection}{ESD CO\_UC\_L7\_1 – getIncomingRequests}
\subsubsection*{ESD CO\_UC\_L7\_1 – getIncomingRequests}

\begin{lstlisting}[basicstyle=\ttfamily\small, breaklines=true]
sequenceDiagram
    autonumber
    actor Lagerhalter
    participant ctrl as :StorageOwnerController
    participant so as so:StorageOwner

    Lagerhalter->>ctrl: getIncomingRequests(storageOwner)
    ctrl->>so: so.getIncomingRequests()
    so-->>ctrl: requests
    ctrl-->>Lagerhalter: Liste von Stellplatzanfragen
\end{lstlisting}

\newpage

% ------------------------------------------------------------------------------
% ESD CO_UC_L7_2 – createOffer
% ------------------------------------------------------------------------------
\phantomsection
\addcontentsline{toc}{subsubsection}{ESD CO\_UC\_L7\_2 – createOffer}
\subsubsection*{ESD CO\_UC\_L7\_2 – createOffer}

\begin{lstlisting}[basicstyle=\ttfamily\small, breaklines=true]
sequenceDiagram
    autonumber
    actor Lagerhalter
    participant ctrl as :StorageOwnerController
    participant so as so:StorageOwner
    participant r as r:Request
    participant o as o:Offer

    Lagerhalter->>ctrl: createOffer(requestId, description, costs)
    ctrl->>so: so.getIncomingRequests()
    so-->>ctrl: requests
    note over ctrl: wählt r mit id = requestId
    ctrl->>so: so.makeOffer(r, description, costs)
    so->>o: Offer(so, r, description, costs)
    o->>o: status = PENDING
    o->>r: r.addOffer(o)
    r->>r: offer = o
    r->>r: status = CLOSED
    so->>so: offers.add(o)
\end{lstlisting}

\newpage

% ------------------------------------------------------------------------------
% ESD CO_UC_L7_3 – getAllOffers
% ------------------------------------------------------------------------------
\phantomsection
\addcontentsline{toc}{subsubsection}{ESD CO\_UC\_L7\_3 – getAllOffers}
\subsubsection*{ESD CO\_UC\_L7\_3 – getAllOffers}

\begin{lstlisting}[basicstyle=\ttfamily\small, breaklines=true]
sequenceDiagram
    autonumber
    actor Lagerhalter
    participant ctrl as :StorageOwnerController
    participant so as so:StorageOwner

    Lagerhalter->>ctrl: getAllOffers(storageOwner)
    ctrl->>so: so.getOffers()
    so-->>ctrl: offers
    ctrl-->>Lagerhalter: Liste von Stellplatzangeboten
\end{lstlisting}

\newpage

% ------------------------------------------------------------------------------
% ESD CO_UC_L7_4 – deleteOffer
% ------------------------------------------------------------------------------
\phantomsection
\addcontentsline{toc}{subsubsection}{ESD CO\_UC\_L7\_4 – deleteOffer}
\subsubsection*{ESD CO\_UC\_L7\_4 – deleteOffer}

\begin{lstlisting}[basicstyle=\ttfamily\small, breaklines=true]
sequenceDiagram
    autonumber
    actor Lagerhalter
    participant ctrl as :StorageOwnerController
    participant so as so:StorageOwner
    participant o as o:Offer

    Lagerhalter->>ctrl: deleteOffer(offerId)
    ctrl->>so: so.getOffers()
    so-->>ctrl: offers
    note over ctrl: wählt o mit id = offerId
    ctrl->>so: so.removeOffer(o)
    so->>so: offers.remove(o)
\end{lstlisting}

\newpage

% ------------------------------------------------------------------------------
% ESD CO_UC_E1_1 – searchByCity
% ------------------------------------------------------------------------------
\phantomsection
\addcontentsline{toc}{subsubsection}{ESD CO\_UC\_E1\_1 – searchByCity}
\subsubsection*{ESD CO\_UC\_E1\_1 – searchByCity}

\begin{lstlisting}[basicstyle=\ttfamily\small, breaklines=true]
sequenceDiagram
    autonumber
    actor Einlagerer
    participant ctrl as :CarOwnerController

    Einlagerer->>ctrl: searchByCity(city)
    ctrl-->>Einlagerer: Liste von Lagerhaltern
\end{lstlisting}

\newpage

% ------------------------------------------------------------------------------
% ESD CO_UC_E2_1 – createRequest
% ------------------------------------------------------------------------------
\phantomsection
\addcontentsline{toc}{subsubsection}{ESD CO\_UC\_E2\_1 – createRequest}
\subsubsection*{ESD CO\_UC\_E2\_1 – createRequest}

\begin{lstlisting}[basicstyle=\ttfamily\small, breaklines=true]
sequenceDiagram
    autonumber
    actor Einlagerer
    participant ctrl as :CarOwnerController
    participant co as co:CarOwner
    participant so as so:StorageOwner
    participant r as r:Request

    Einlagerer->>ctrl: createRequest(storageOwnerId, description)
    ctrl->>co: co.submitRequest(so, description)
    co->>r: Request(co, so, description)
    r->>r: type = STELLPLATZ
    r->>r: status = OPEN
    co->>co: requests.add(r)
    co->>so: so.addRequest(r)
\end{lstlisting}

\newpage

% ------------------------------------------------------------------------------
% ESD CO_UC_E2_2 – getAllRequests
% ------------------------------------------------------------------------------
\phantomsection
\addcontentsline{toc}{subsubsection}{ESD CO\_UC\_E2\_2 – getAllRequests}
\subsubsection*{ESD CO\_UC\_E2\_2 – getAllRequests}

\begin{lstlisting}[basicstyle=\ttfamily\small, breaklines=true]
sequenceDiagram
    autonumber
    actor Einlagerer
    participant ctrl as :CarOwnerController
    participant co as co:CarOwner

    Einlagerer->>ctrl: getAllRequests(carOwner)
    ctrl->>co: co.getRequests()
    co-->>ctrl: requests
    ctrl-->>Einlagerer: Liste von Stellplatzanfragen
\end{lstlisting}

\newpage

% ------------------------------------------------------------------------------
% ESD CO_UC_E2_3 – deleteRequest
% ------------------------------------------------------------------------------
\phantomsection
\addcontentsline{toc}{subsubsection}{ESD CO\_UC\_E2\_3 – deleteRequest}
\subsubsection*{ESD CO\_UC\_E2\_3 – deleteRequest}

\begin{lstlisting}[basicstyle=\ttfamily\small, breaklines=true]
sequenceDiagram
    autonumber
    actor Einlagerer
    participant ctrl as :CarOwnerController
    participant co as co:CarOwner

    Einlagerer->>ctrl: deleteRequest(requestId)
    ctrl->>co: co.getRequests()
    co-->>ctrl: requests
    note over ctrl: wählt r mit id = requestId
    ctrl->>co: co.removeRequest(r)
    co->>co: requests.remove(r)
\end{lstlisting}

\newpage

% ------------------------------------------------------------------------------
% ESD CO_UC_E3_1 – getReceivedOffers
% ------------------------------------------------------------------------------
\phantomsection
\addcontentsline{toc}{subsubsection}{ESD CO\_UC\_E3\_1 – getReceivedOffers}
\subsubsection*{ESD CO\_UC\_E3\_1 – getReceivedOffers}

\begin{lstlisting}[basicstyle=\ttfamily\small, breaklines=true]
sequenceDiagram
    autonumber
    actor Einlagerer
    participant ctrl as :CarOwnerController
    participant co as co:CarOwner
    participant r as r:Request

    Einlagerer->>ctrl: getReceivedOffers(carOwner)
    ctrl->>co: co.getRequests()
    co-->>ctrl: requests

    loop Fuer jeden Request in requests
        ctrl->>r: r.getOffer()
        r-->>ctrl: offer
    end

    ctrl-->>Einlagerer: Liste von Stellplatzangeboten
\end{lstlisting}

\newpage

% ------------------------------------------------------------------------------
% ESD CO_UC_E3_2 – respondToOffer
% ------------------------------------------------------------------------------
\phantomsection
\addcontentsline{toc}{subsubsection}{ESD CO\_UC\_E3\_2 – respondToOffer}
\subsubsection*{ESD CO\_UC\_E3\_2 – respondToOffer}

\begin{lstlisting}[basicstyle=\ttfamily\small, breaklines=true]
sequenceDiagram
    autonumber
    actor Einlagerer
    participant ctrl as :CarOwnerController
    participant co as co:CarOwner
    participant o as o:Offer

    Einlagerer->>ctrl: respondToOffer(offerId, accepted)
    ctrl->>co: co.getRequests()
    co-->>ctrl: requests
    note over ctrl: wählt o mit id = offerId

    alt accepted = true
        ctrl->>o: o.accept()
        o->>o: status = ACCEPTED
    else accepted = false
        ctrl->>o: o.reject()
        o->>o: status = REJECTED
    end
\end{lstlisting}

\newpage

% ------------------------------------------------------------------------------
% ESD CO_UC_E5_1 – createAppointment
% ------------------------------------------------------------------------------
\phantomsection
\addcontentsline{toc}{subsubsection}{ESD CO\_UC\_E5\_1 – createAppointment}
\subsubsection*{ESD CO\_UC\_E5\_1 – createAppointment}

\begin{lstlisting}[basicstyle=\ttfamily\small, breaklines=true]
sequenceDiagram
    autonumber
    actor Lagerhalter
    participant ctrl as :StorageOwnerController
    participant so as so:StorageOwner
    participant a as a:Appointment

    Lagerhalter->>ctrl: createAppointment(type, date, time)
    ctrl->>so: so.scheduleAppointment(type, date, time)
    so->>a: Appointment(so, type, date, time)
    a->>a: subject = null
    a->>a: status = AVAILABLE
    so->>so: appointments.add(a)
\end{lstlisting}

\newpage

% ------------------------------------------------------------------------------
% ESD CO_UC_E5_2 – getAvailableAppointments
% ------------------------------------------------------------------------------
\phantomsection
\addcontentsline{toc}{subsubsection}{ESD CO\_UC\_E5\_2 – getAvailableAppointments}
\subsubsection*{ESD CO\_UC\_E5\_2 – getAvailableAppointments}

\begin{lstlisting}[basicstyle=\ttfamily\small, breaklines=true]
sequenceDiagram
    autonumber
    actor Einlagerer
    participant ctrl as :CarOwnerController
    participant so as so:StorageOwner

    Einlagerer->>ctrl: getAvailableAppointments(carOwnerId, storageOwnerId)
    ctrl->>so: so.getAppointments()
    so-->>ctrl: appointments
    note over ctrl: filtert Termine mit status = AVAILABLE
    ctrl-->>Einlagerer: Liste von Terminen
\end{lstlisting}

\newpage

% ------------------------------------------------------------------------------
% ESD CO_UC_E5_3 – bookAppointment
% ------------------------------------------------------------------------------
\phantomsection
\addcontentsline{toc}{subsubsection}{ESD CO\_UC\_E5\_3 – bookAppointment}
\subsubsection*{ESD CO\_UC\_E5\_3 – bookAppointment}

\begin{lstlisting}[basicstyle=\ttfamily\small, breaklines=true]
sequenceDiagram
    autonumber
    actor Einlagerer
    participant ctrl as :CarOwnerController
    participant co as co:CarOwner
    participant so as so:StorageOwner
    participant a as a:Appointment
    participant c as c:Car

    Einlagerer->>ctrl: bookAppointment(storageOwnerId, carId, type, date, time)
    note over ctrl: wählt Fahrzeug c und verfügbaren Termin a
    ctrl->>c: c.relocateTo(newStorageOwner)
    alt type = EINLAGERUNG
        c->>so: so.addCarToStorage(c)
        so->>c: c.setCurrentStorage(so)
    else type = AUSLAGERUNG
        c->>so: so.removeCarFromStorage(c)
        so->>c: c.setCurrentStorage(null)
    end
    ctrl->>a: a.book(co, c)
    a->>a: status = BOOKED
    a->>a: carOwner = co
    a->>a: subject = c
    ctrl->>co: co.addAppointment(a)
\end{lstlisting}

\newpage

% ------------------------------------------------------------------------------
% ESD CO_UC_E5_4 – getAllBookedAppointments
% ------------------------------------------------------------------------------
\phantomsection
\addcontentsline{toc}{subsubsection}{ESD CO\_UC\_E5\_4 – getAllBookedAppointments}
\subsubsection*{ESD CO\_UC\_E5\_4 – getAllBookedAppointments}

\begin{lstlisting}[basicstyle=\ttfamily\small, breaklines=true]
sequenceDiagram
    autonumber
    actor Einlagerer
    participant ctrl as :CarOwnerController
    participant co as co:CarOwner

    Einlagerer->>ctrl: getAllBookedAppointments(carOwner)
    ctrl->>co: co.getBookedAppointments()
    co-->>ctrl: bookedAppointments
    ctrl-->>Einlagerer: Liste von Terminen
\end{lstlisting}

\newpage

% ------------------------------------------------------------------------------
% ESD CO_UC_E5_5 – cancelAppointment
% ------------------------------------------------------------------------------
\phantomsection
\addcontentsline{toc}{subsubsection}{ESD CO\_UC\_E5\_5 – cancelAppointment}
\subsubsection*{ESD CO\_UC\_E5\_5 – cancelAppointment}

\begin{lstlisting}[basicstyle=\ttfamily\small, breaklines=true]
sequenceDiagram
    autonumber
    actor Einlagerer
    participant ctrl as :CarOwnerController
    participant co as co:CarOwner
    participant a as a:Appointment

    Einlagerer->>ctrl: cancelAppointment(appointmentId)
    ctrl->>co: co.getBookedAppointments()
    co-->>ctrl: bookedAppointments
    note over ctrl: wählt a mit id = appointmentId
    ctrl->>a: a.cancel()
    a->>a: status = CANCELLED
\end{lstlisting}

\newpage

% ------------------------------------------------------------------------------
% ESD CO_UC_E5_6 – getAllAppointments
% ------------------------------------------------------------------------------
\phantomsection
\addcontentsline{toc}{subsubsection}{ESD CO\_UC\_E5\_6 – getAllAppointments}
\subsubsection*{ESD CO\_UC\_E5\_6 – getAllAppointments}

\begin{lstlisting}[basicstyle=\ttfamily\small, breaklines=true]
sequenceDiagram
    autonumber
    actor Lagerhalter
    participant ctrl as :StorageOwnerController
    participant so as so:StorageOwner

    Lagerhalter->>ctrl: getAllAppointments(storageOwner)
    ctrl->>so: so.getAppointments()
    so-->>ctrl: appointments
    ctrl-->>Lagerhalter: Liste von Terminen
\end{lstlisting}

\newpage

% ------------------------------------------------------------------------------
% ESD CO_UC_E5_7 – deleteAppointment
% ------------------------------------------------------------------------------
\phantomsection
\addcontentsline{toc}{subsubsection}{ESD CO\_UC\_E5\_7 – deleteAppointment}
\subsubsection*{ESD CO\_UC\_E5\_7 – deleteAppointment}

\begin{lstlisting}[basicstyle=\ttfamily\small, breaklines=true]
sequenceDiagram
    autonumber
    actor Lagerhalter
    participant ctrl as :StorageOwnerController
    participant so as so:StorageOwner
    participant a as a:Appointment

    Lagerhalter->>ctrl: deleteAppointment(appointmentId)
    ctrl->>so: so.getAppointments()
    so-->>ctrl: appointments
    note over ctrl: wählt a mit id = appointmentId
    ctrl->>so: so.removeAppointment(a)
    so->>so: appointments.remove(a)
\end{lstlisting}

\newpage

% ------------------------------------------------------------------------------
% ESD CO_UC_E6_1 – bookOptionalService
% ------------------------------------------------------------------------------
\phantomsection
\addcontentsline{toc}{subsubsection}{ESD CO\_UC\_E6\_1 – bookOptionalService}
\subsubsection*{ESD CO\_UC\_E6\_1 – bookOptionalService}

\begin{lstlisting}[basicstyle=\ttfamily\small, breaklines=true]
sequenceDiagram
    autonumber
    actor Einlagerer
    participant ctrl as :CarOwnerController
    participant c as c:Car
    participant s as s:Service

    Einlagerer->>ctrl: bookOptionalService(carId, storageOwnerId, serviceId)
    note over ctrl: wählt Fahrzeug c und Service s
    ctrl->>c: c.addBookedOptionalService(s)
    c->>c: bookedServices.add(s)
\end{lstlisting}

\newpage

% ------------------------------------------------------------------------------
% ESD CO_UC_E6_2 – removeBookedService
% ------------------------------------------------------------------------------
\phantomsection
\addcontentsline{toc}{subsubsection}{ESD CO\_UC\_E6\_2 – removeBookedService}
\subsubsection*{ESD CO\_UC\_E6\_2 – removeBookedService}

\begin{lstlisting}[basicstyle=\ttfamily\small, breaklines=true]
sequenceDiagram
    autonumber
    actor Einlagerer
    participant ctrl as :CarOwnerController
    participant c as c:Car

    Einlagerer->>ctrl: removeBookedService(carId, serviceId)
    note over ctrl: wählt Fahrzeug c und gebuchten Service s
    ctrl->>c: c.removeBookedOptionalService(s)
    c->>c: bookedServices.remove(s)
\end{lstlisting}

\newpage

% ------------------------------------------------------------------------------
% ESD CO_UC_E7_1 – createCar
% ------------------------------------------------------------------------------
\phantomsection
\addcontentsline{toc}{subsubsection}{ESD CO\_UC\_E7\_1 – createCar}
\subsubsection*{ESD CO\_UC\_E7\_1 – createCar}

\begin{lstlisting}[basicstyle=\ttfamily\small, breaklines=true]
sequenceDiagram
    autonumber
    actor Einlagerer
    participant ctrl as :CarOwnerController
    participant co as co:CarOwner
    participant c as c:Car

    Einlagerer->>ctrl: createCar(licensePlate)
    ctrl->>co: co.registerCar(licensePlate)
    co->>c: Car(co, licensePlate)
    c->>c: readyToDrive = true
    c->>c: isMaintained = false
    co->>co: cars.add(c)
\end{lstlisting}

\newpage

% ------------------------------------------------------------------------------
% ESD CO_UC_E8_1 – getAllCars
% ------------------------------------------------------------------------------
\phantomsection
\addcontentsline{toc}{subsubsection}{ESD CO\_UC\_E8\_1 – getAllCars}
\subsubsection*{ESD CO\_UC\_E8\_1 – getAllCars}

\begin{lstlisting}[basicstyle=\ttfamily\small, breaklines=true]
sequenceDiagram
    autonumber
    actor Einlagerer
    participant ctrl as :CarOwnerController
    participant co as co:CarOwner

    Einlagerer->>ctrl: getAllCars(carOwner)
    ctrl->>co: co.getCars()
    co-->>ctrl: cars
    ctrl-->>Einlagerer: Liste von Fahrzeugen
\end{lstlisting}

\newpage

% ------------------------------------------------------------------------------
% ESD CO_UC_E8_2 – updateCar
% ------------------------------------------------------------------------------
\phantomsection
\addcontentsline{toc}{subsubsection}{ESD CO\_UC\_E8\_2 – updateCar}
\subsubsection*{ESD CO\_UC\_E8\_2 – updateCar}

\begin{lstlisting}[basicstyle=\ttfamily\small, breaklines=true]
sequenceDiagram
    autonumber
    actor Einlagerer
    participant ctrl as :CarOwnerController
    participant co as co:CarOwner
    participant c as c:Car

    Einlagerer->>ctrl: updateCar(carId, licensePlate)
    ctrl->>co: co.getCars()
    co-->>ctrl: cars
    note over ctrl: wählt c mit id = carId
    ctrl->>c: c.changeLicensePlate(licensePlate)
\end{lstlisting}

\newpage

% ------------------------------------------------------------------------------
% ESD CO_UC_E8_3 – removeCar
% ------------------------------------------------------------------------------
\phantomsection
\addcontentsline{toc}{subsubsection}{ESD CO\_UC\_E8\_3 – removeCar}
\subsubsection*{ESD CO\_UC\_E8\_3 – removeCar}

\begin{lstlisting}[basicstyle=\ttfamily\small, breaklines=true]
sequenceDiagram
    autonumber
    actor Einlagerer
    participant ctrl as :CarOwnerController
    participant co as co:CarOwner

    Einlagerer->>ctrl: removeCar(carId)
    ctrl->>co: co.getCars()
    co-->>ctrl: cars
    note over ctrl: wählt c mit id = carId
    ctrl->>co: co.removeCar(c)
    co->>co: cars.remove(c)
\end{lstlisting}

\newpage
